% !TeX program = lualatex
% ============================================================
% chapter8.tex — ch08 唯一內容源（2026-08-09 起；LaTeX 統一 U1）
% 沿革：升格自 dist/ch08/chapter8.tex（P1 首轉：make_dist.py 自 fragment 轉換；
%   fragment 樹已凍結，改課文一律改本檔。拍板見 ../../KICKOFF-latex-unification.md）
%   樣式源＝../../template/calcbook.sty（M-B1 拍板模板）
% 編譯（本資料夾為 CWD；aux 進 build/、成品 PDF 移入 ../../dist/ch08/）：
%   latexmk -lualatex -auxdir=../../build/aux-ch08 chapter8.tex
% ============================================================
\documentclass[a4paper,12pt,oneside]{memoir}
\makeatletter\def\input@path{{../../template/}}\makeatother
\def\cbfontsdir{../../template/fonts/inter/}
\usepackage{calcbook}
% 圖：emitter 射 `<ch>/<stem>`（見 convert.py），故 graphicspath 指向 chapters/ 上層。
% 模板內的 {../figs/} 是 chapters/<ch>/ 為 CWD 時的路徑，dist/<ch>/ 需這一行覆寫。
\graphicspath{{../../chapters/}}
\begin{document}
\cbchapter{8}
\chapteropener{Chapter 8}{Techniques of Integration}
\begin{lead}
The Fundamental Theorem of Calculus reduced the evaluation of a definite integral to a single task: find an antiderivative. Chapter 6 supplied one tool for that task, substitution, and Chapter 7 showed how much geometry and physics the integral can carry. But substitution alone reaches only integrands with a visible inner function. Products such as \(x\cos x\), powers such as \(\sin^4 x\), and rational functions such as \(\tfrac{x+1}{x(x^2+1)}\) all resist it, and §6.5 promised that their techniques would arrive here. This chapter completes the symbolic toolkit: integration by parts, trigonometric methods, and partial fractions, followed by training in how to choose among them. The chapter then widens the integral itself to unbounded intervals and unbounded integrands, and closes with numerical rules, with error control, for the integrals that no formula reaches.
\end{lead}
\parahead{By the end of this chapter, you will be able to:}
\begin{objectives}
  \item integrate products, logarithms, and inverse trigonometric functions by parts;
  \item evaluate trigonometric integrals with the standard identities, and integrals containing quadratic radicals by trigonometric substitution;
  \item integrate rational functions by partial fractions, and choose an approach for an unfamiliar integral;
  \item decide whether an improper integral converges, by evaluation or by comparison, and evaluate it when it does;
  \item approximate a definite integral with the Midpoint, Trapezoidal, and Simpson’s rules, and bound the error.
\end{objectives}
\sechead{8.1}{Integration by Parts}
Try to compute \(\int x\cos x\,dx\) by substitution. The method needs two things, an inner function and its derivative alongside, and this integrand has neither: \(x\) and \(\cos x\) are unrelated factors, and no choice of \(u\) makes the rest of the integrand equal \(du\). Substitution reverses the chain rule, so it applies when the integrand looks like the output of the chain rule. The integrand \(x \cos x\) does not. It looks instead like one term in the output of the \emph{product rule}, a rule we have not yet run backwards. Running it backwards gives the second technique of integration, and the subject of this section.

Start from the product rule itself. If \(f\) and \(g\) are differentiable, then by Theorem 2.6,

\[ \bigl[f(x)\,g(x)\bigr]' = f'(x)\,g(x) + f(x)\,g'(x). \]

Read this as a statement about antiderivatives (Definition 6.3): it says that \(fg\) is an antiderivative of the sum \(f'g + fg'\). Solving for the term we want turns the rule into a tool.

\begin{envtheorem}{Theorem}{thm:8.1}{Integration by Parts}
Let \(f\) and \(g\) be differentiable on an open interval \(I\). Then, on \(I\),

\[ \int f(x)\,g'(x)\,dx = f(x)\,g(x) - \int f'(x)\,g(x)\,dx. \]

Precisely: if \(G\) is any antiderivative of \(f'g\) on \(I\), then \(fg - G\) is an antiderivative of \(fg'\) on \(I\).
\end{envtheorem}
\begin{envproof}{Proof}{}{}
Differentiate \(fg - G\). By the product rule (Theorem 2.6) and the definition of \(G\),

\[ (fg - G)' = \bigl(f'g + fg'\bigr) - f'g = fg'. \]

So \(fg - G\) is an antiderivative of \(fg'\), which is what the displayed identity asserts. \qedmark
\end{envproof}
An equation between indefinite integrals is an equation between families of functions. The constant of integration is shared: it appears once, when the remaining integral on the right is finally evaluated. The identity trades one integral for another. That is all it does, and that is its value: the trade is progress whenever the new integrand \(f'g\) is easier than the old one \(fg'\).

In practice the formula is applied through the differential notation of §5.3, in the same way substitution was in §6.5. Set \(u = f(x)\) and \(v = g(x)\), so that \(du = f'(x)\,dx\) and \(dv = g'(x)\,dx\). The formula then reads:

\[ \int u\,dv = uv - \int v\,du. \]

To use it, split the given integrand into two parts, a function \(u\) and a differential \(dv\); the name of the technique records that split. From \(u\) we compute \(du\) by differentiating, and from \(dv\) we compute \(v\) by antidifferentiating. Any antiderivative serves for \(v\), so we take the one with zero constant.

Dropping the constant changes nothing: replacing \(v\) by \(v + k\) adds \(ku\) to \(uv\) and adds the same \(ku\) to \(\int v\,du\), which enters with a minus sign, so the two changes cancel. Occasionally a different choice of \(v\) is worth making, because it can simplify the remaining integral. In \(\int\ln(x + 1)\,dx\), taking \(v = x + 1\) instead of \(v = x\) leaves \(\int\frac{x+1}{x+1}\,dx = \int dx\), while \(v = x\) leaves \(\int\frac{x}{x+1}\,dx\), which needs a division step first. Both routes reach the same answer.

\begin{figureblock}
  \includegraphics[width=79.89mm]{ch08/figs/parts-uv-rectangle}
\figcaption{fig:8.1}{The parts formula as area bookkeeping. For a curve in the \(uv\)-plane climbing from the origin to \((u_{1}, v_{1})\) with \(u\) and \(v\) increasing together, the rectangle of area \(u_{1}v_{1}\) splits into \(\int v\,du\), below the curve, and \(\int u\,dv\), beside it; subtracting either from the rectangle leaves the other. The picture illustrates the formula; the proof remains that of Theorem \ref{thm:8.1}.}
\end{figureblock}
\begin{envstrategy}{Strategy}{strat:8.1}{Choosing \(u\) and \(dv\)}
\begin{steps}
  \item \emph{Pick \(u\) to get simpler when differentiated.} A power of \(x\) drops degree; a logarithm or inverse trigonometric function turns algebraic. Pick \(dv\) as the rest of the integrand, and check that you can antidifferentiate \(dv\) on its own.
  \item \emph{Logarithms and inverse trigonometric functions go into \(u\).} Their derivatives are algebraic, and they have no elementary \(dv\)-role: taking \(dv = \ln x\,dx\) would require the very antiderivative we are seeking. When the integrand is nothing but \(\ln x\) or \(\arctan x\), take \(dv = dx\) (Examples \ref{ex:8.2} and \ref{ex:8.5}).
  \item \emph{If the integrand reproduces itself, solve for the integral.} When two applications of parts return the original integral, as with \(e^{x}\sin x\), collect that integral on one side of the equation and divide (Example \ref{ex:8.4}).
  \item \emph{If every choice makes the integral worse, parts is the wrong tool.} Return to substitution, or to the methods of §8.2 through §8.4.
\end{steps}
\end{envstrategy}
\begin{workedexample}
\begin{envexample}{Example}{ex:8.1}{}
Find \(\displaystyle\int x\cos x\,dx\).
\end{envexample}
\begin{envsolution}{Solution}{}{}
Following step 1 of Strategy \ref{strat:8.1}, take \(u = x\), since differentiation lowers its degree, and \(dv = \cos x\,dx\), which we can antidifferentiate. Then \(du = dx\) and \(v = \sin x\), and the formula gives

\[ \int x\cos x\,dx = x\sin x - \int \sin x\,dx = x\sin x + \cos x + C. \]

Differentiating the answer confirms it: \((x\sin x + \cos x)' = \sin x + x\cos x - \sin x = x\cos x\).

The choice of parts matters. Had we taken \(u = \cos x\) and \(dv = x\,dx\), so that \(du = -\sin x\,dx\) and \(v = \tfrac{x^{2}}{2}\), the trade would read

\[ \int x\cos x\,dx = \frac{x^{2}}{2}\cos x + \frac{1}{2}\int x^{2}\sin x\,dx, \]

which is true but worse: the power of \(x\) went up. A trade should simplify, and step 1 is the test to apply before trading. \qedmark
\end{envsolution}
\end{workedexample}
\begin{workedexample}
\begin{envexample}{Example}{ex:8.2}{}
Find \(\displaystyle\int \ln x\,dx\).
\end{envexample}
\begin{envsolution}{Solution}{}{}
The integrand does not look like a product, but step 2 of Strategy \ref{strat:8.1} supplies the split: take \(u = \ln x\) and let \(dv = dx\) be the invisible factor. Then \(du = \tfrac{1}{x}\,dx\) by Theorem 4.14, and \(v = x\). For \(x > 0\), where \(\ln x\) is defined,

\[ \int \ln x\,dx = x\ln x - \int x\cdot\frac{1}{x}\,dx = x\ln x - \int 1\,dx = x\ln x - x + C. \]

The trade turned a logarithm into the constant integrand \(1\). This is the pattern of step 2 in action: placing the logarithm in \(u\) produced the algebraic derivative \(\tfrac{1}{x}\), and \(dv\) was nothing more than \(dx\). \qedmark
\end{envsolution}
\end{workedexample}
\begin{workedexample}
\begin{envexample}{Example}{ex:8.3}{}
Find \(\displaystyle\int t^{2}e^{t}\,dt\).
\end{envexample}
\begin{envsolution}{Solution}{}{}
Take \(u = t^{2}\), whose degree drops under differentiation, and \(dv = e^{t}\,dt\), so \(du = 2t\,dt\) and \(v = e^{t}\):

\[ \int t^{2}e^{t}\,dt = t^{2}e^{t} - 2\int t e^{t}\,dt. \]

The degree dropped from \(2\) to \(1\), and one more round of parts finishes the job. In the remaining integral take \(u = t\), \(dv = e^{t}\,dt\), so \(du = dt\), \(v = e^{t}\):

\[ \int t e^{t}\,dt = t e^{t} - \int e^{t}\,dt = t e^{t} - e^{t} + C_1. \]

Substituting back, with the constant \(-2C_{1}\) absorbed into \(C\),

\[ \int t^{2}e^{t}\,dt = t^{2}e^{t} - 2\bigl(t e^{t} - e^{t}\bigr) + C = \bigl(t^{2} - 2t + 2\bigr)e^{t} + C. \]

Each application of parts lowers the power by one, so an integrand \(t^{n}e^{t}\) yields to \(n\) rounds. \qedmark
\end{envsolution}
\end{workedexample}
\begin{workedexample}
\begin{envexample}{Example}{ex:8.4}{}
Find \(\displaystyle\int e^{x}\sin x\,dx\).
\end{envexample}
\begin{envsolution}{Solution}{}{}
Neither factor gets simpler under differentiation: \(e^{x}\) reproduces itself and \(\sin x\) cycles. Parts will not simplify this integral, but by step 3 of Strategy \ref{strat:8.1} it can still close it. Write \(J = \int e^{x}\sin x\,dx\) for the unknown integral. Take \(u = \sin x\), \(dv = e^{x}\,dx\), so \(du = \cos x\,dx\), \(v = e^{x}\):

\[ J = e^{x}\sin x - \int e^{x}\cos x\,dx. \]

Apply parts once more to the new integral, with \(u = \cos x\), \(dv = e^{x}\,dx\), so \(du = -\sin x\,dx\), \(v = e^{x}\):

\[ \int e^{x}\cos x\,dx = e^{x}\cos x + \int e^{x}\sin x\,dx = e^{x}\cos x + J. \]

Substituting back,

\[ J = e^{x}\sin x - e^{x}\cos x - J, \]

and solving for \(J\) gives \(2J = e^{x}(\sin x - \cos x)\), so

\[ \int e^{x}\sin x\,dx = \frac{e^{x}\bigl(\sin x - \cos x\bigr)}{2} + C. \]

Here we treated the unknown integral as a single quantity and solved for it. The constant of integration is restored at the end, since the equation holds for antiderivative families. Differentiating the answer confirms it: \(\tfrac{1}{2}\bigl[e^{x}(\sin x - \cos x) + e^{x}(\cos x + \sin x)\bigr] = e^{x}\sin x\). \qedmark
\end{envsolution}
\end{workedexample}
\begin{envcaution}{Caution}{}{}
Solving for the integral depends on keeping the same kind of factor in the \(u\) position at every round. Reversing the roles on the second application undoes the first. In Example \ref{ex:8.4}, choosing \(u = e^{x}\) and \(dv = \cos x\,dx\) at the second step gives \(\int e^{x}\cos x\,dx = e^{x}\sin x - J\), and substituting that back into the first equation returns \(J = J\), a true statement that carries no information. When a second application of parts yields an identity instead of an equation for the unknown, this is why.
\end{envcaution}
For definite integrals, the term \(uv\) is evaluated between the limits: it becomes \(\bigl[f(x)\,g(x)\bigr]_{a}^{b}\) in the next theorem. The hypotheses match those of the Fundamental Theorem: we use the term \emph{smooth on \([a, b]\)} from §7.6, meaning that \(f\) is differentiable with \(f'\) continuous on \([a, b]\), one-sided at the endpoints.

\begin{envtheorem}{Theorem}{thm:8.2}{Integration by Parts for Definite Integrals}
Let \(f\) and \(g\) be smooth on \([a, b]\). Then

\[ \int_{a}^{b} f(x)\,g'(x)\,dx = \Bigl[f(x)\,g(x)\Bigr]_{a}^{b} - \int_{a}^{b} f'(x)\,g(x)\,dx. \]
\end{envtheorem}
\begin{envproof}{Proof}{}{}
The sum \(f'g + fg'\) is continuous on \([a, b]\), since \(f, g, f', g'\) all are. Even with these hypotheses, Theorem 6.4 cannot be cited directly, because it asks for an antiderivative on an open interval larger than \([a, b]\), and smoothness on \([a, b]\) alone does not provide one; instead we build the antiderivative we need. Let \(H(x) = \int_{a}^{x} \bigl[f'(t)\,g(t) + f(t)\,g'(t)\bigr]\,dt\). By the Fundamental Theorem of Calculus, Part 1 (Theorem 6.3), \(H\) is continuous on \([a, b]\) with \(H' = f'g + fg'\) on \((a, b)\). The product \(fg\) has the same derivative there, by the product rule (Theorem 2.6). The difference \(fg - H\) is therefore continuous on \([a, b]\) with zero derivative on \((a, b)\), hence constant (Corollary 4.4). Evaluating at \(b\) and at \(a\), where \(H(a) = 0\), gives

\[ \int_{a}^{b} \bigl[f'(x)\,g(x) + f(x)\,g'(x)\bigr]\,dx = f(b)\,g(b) - f(a)\,g(a). \]

Splitting the left side by linearity (Theorem 6.2, part 2) and moving \(\int_{a}^{b} f'g\) to the right side gives the formula. \qedmark
\end{envproof}
\begin{workedexample}
\begin{envexample}{Example}{ex:8.5}{}
Evaluate \(\displaystyle\int_{0}^{1} \arctan x\,dx\).
\end{envexample}
\begin{envsolution}{Solution}{}{}
As in Example \ref{ex:8.2}, the integrand is a lone function whose derivative is algebraic, so take \(u = \arctan x\) and \(dv = dx\). Then \(du = \tfrac{1}{1+x^{2}}\,dx\), by Example 3.16, and \(v = x\). Both \(\arctan x\) and \(x\) are smooth on \([0, 1]\), so Theorem \ref{thm:8.2} applies:

\[ \int_{0}^{1} \arctan x\,dx = \Bigl[x\arctan x\Bigr]_{0}^{1} - \int_{0}^{1} \frac{x}{1+x^{2}}\,dx = \frac{\pi}{4} - \int_{0}^{1} \frac{x}{1+x^{2}}\,dx, \]

since \(\arctan 1 = \tfrac{\pi}{4}\) and the \(x = 0\) term vanishes. The remaining integral is a substitution: with \(w = 1 + x^{2}\), \(dw = 2x\,dx\), and the limits transform by Theorem 6.7 from \(x \in [0, 1]\) to \(w \in [1, 2]\):

\[ \int_{0}^{1} \frac{x}{1+x^{2}}\,dx = \frac{1}{2}\int_{1}^{2} \frac{dw}{w} = \frac{1}{2}\Bigl[\ln w\Bigr]_{1}^{2} = \frac{\ln 2}{2}. \]

Therefore

\[ \int_{0}^{1} \arctan x\,dx = \frac{\pi}{4} - \frac{\ln 2}{2}. \]

One integral, two techniques: parts opened it, substitution finished it. That division of labor is common, and §8.5 trains the choices. \qedmark
\end{envsolution}
\end{workedexample}
\begin{envcaution}{Caution}{}{}
The minus sign in \(uv - \int v\,du\) applies to the \emph{entire} remaining integral. When \(\int v\,du\) is itself evaluated by another round of parts, as in Examples \ref{ex:8.3} and \ref{ex:8.4}, enclose its whole result in parentheses before distributing the sign. Dropped parentheses at this step are the most common error in multi-round parts.
\end{envcaution}
\subsechead{Reduction formulas}
Parts can lower an exponent by a fixed step even when it cannot remove it outright. Applied once to \(\int\sin^{n}x\,dx\), the formula relates the integral to the integral two exponents down. The result is worth recording once, in general, so that any particular power is a chain of lookups rather than a fresh computation.

\begin{workedexample}
\begin{envexample}{Example}{ex:8.6}{}
Let \(n \ge 2\) be an integer. Show that

\[ \int \sin^{n}x\,dx = -\frac{1}{n}\cos x\,\sin^{n-1}x + \frac{n-1}{n}\int \sin^{n-2}x\,dx. \]
\end{envexample}
\begin{envsolution}{Solution}{}{}
Split off one factor of \(\sin x\) to serve as \(dv\): take \(u = \sin^{n-1}x\) and \(dv = \sin x\,dx\), so that \(du = (n-1)\sin^{n-2}x\cos x\,dx\) by the chain rule and \(v = -\cos x\). Parts gives

\[ \int \sin^{n}x\,dx = -\cos x\,\sin^{n-1}x + (n-1)\int \sin^{n-2}x\,\cos^{2}x\,dx. \]

Now \(\cos^{2}x = 1 - \sin^{2}x\), so the right-hand integral splits:

\[ \begin{aligned}
          \int \sin^{n}x\,dx &= -\cos x\,\sin^{n-1}x + (n-1)\int \sin^{n-2}x\,dx \\
            &\quad - (n-1)\int \sin^{n}x\,dx.
        \end{aligned} \]

The original integral has reappeared on the right, and as in Example \ref{ex:8.4} we solve for it. Adding \((n-1)\int\sin^{n}x\,dx\) to both sides and dividing by \(n\):

\[ \int \sin^{n}x\,dx = -\frac{1}{n}\cos x\,\sin^{n-1}x + \frac{n-1}{n}\int \sin^{n-2}x\,dx. \]

As a check, set \(n = 2\). The formula gives \(\int\sin^{2}x\,dx = -\tfrac{1}{2}\cos x\sin x + \tfrac{1}{2}x + C\), and since \(\sin x\cos x = \tfrac{1}{2}\sin 2x\), this is \(\tfrac{x}{2} - \tfrac{\sin 2x}{4} + C\). That is the same antiderivative that the power-reduction identity \(\sin^{2}x = \tfrac{1 - \cos 2x}{2}\) of §A.2 produces by direct integration. Repeated powers of sine and cosine are the opening business of the next section, which uses both routes. \qedmark
\end{envsolution}
\end{workedexample}
The formula pays off most visibly on definite integrals over \([0, \tfrac{\pi}{2}]\), where the boundary term disappears. Both \(\sin^{n-1}x\) and \(-\cos x\) are smooth there, so Theorem \ref{thm:8.2} applies and the boundary term can be evaluated. It vanishes at both ends:

\[ \Bigl[-\tfrac{1}{n}\cos x\,\sin^{n-1}x\Bigr]_{0}^{\pi/2} = 0, \]

because \(\cos\tfrac{\pi}{2} = 0\) at the upper limit and \(\sin^{n-1}0 = 0\) at the lower one, the second of these because \(n \ge 2\). Writing \(I_{n} = \int_{0}^{\pi/2}\sin^{n}x\,dx\), the reduction formula collapses to a relation between numbers,

\[ I_{n} = \frac{n-1}{n}\,I_{n-2} \qquad (n \ge 2), \]

and every even power now comes from repeated lookup, ending at \(I_{0} = \int_{0}^{\pi/2} 1\,dx = \tfrac{\pi}{2}\). For \(n = 6\),

\[ I_{6} = \tfrac{5}{6}I_{4} = \tfrac{5}{6}\cdot\tfrac{3}{4}\,I_{2} = \tfrac{5}{6}\cdot\tfrac{3}{4}\cdot\tfrac{1}{2}\,I_{0} = \tfrac{5}{6}\cdot\tfrac{3}{4}\cdot\tfrac{1}{2}\cdot\tfrac{\pi}{2} = \frac{5\pi}{32}, \]

and no integration was carried out after the formula itself. That is the value of recording a reduction formula once.

\subsechead{Settling the shell–washer question}
Section 7.3 computed one solid’s volume by shells and by washers, found the answers equal, and recorded a promise: the tools to explain the agreement would arrive with integration by parts in Chapter 8 and with change of variables in Chapter 15. Parts has now arrived, and it settles the question for the solid §7.3 actually compared, and for every solid of that monotone kind.

Let \(f\) be strictly increasing and smooth on \([a, b]\), with \(0 \le a\) and \(f \ge 0\), and suppose its inverse \(g\) is continuous on \([c, d]\), where \(c = f(a)\) and \(d = f(b)\). Rotate the region under \(y = f(x)\) about the \(y\)-axis. In §7.3’s Example 7.12 this is the case \(f(x) = x^{2}\) on \([0, 1]\), with \(g(y) = \sqrt{y}\). The shell method (Theorem 7.1) gives the volume as \(\int_a^b 2\pi x\,f(x)\,dx\).

The washer reading slices horizontally instead. Below height \(c\) every slice is the full annulus from \(x = a\) to \(x = b\). At height \(y\) between \(c\) and \(d\), the region contains exactly the points with \(f(x) \ge y\); since \(f\) is strictly increasing, that condition says exactly that \(x \ge g(y)\), so the slice runs from \(x = g(y)\) out to \(x = b\). Figure \ref{fig:8.2} draws the split. Applying the cross-section definition (Definition 7.2) on each of \([0, c]\) and \([c, d]\) and adding, the washer volume is

\[ V_{\text{washer}} = \pi c\,\bigl(b^{2} - a^{2}\bigr) + \pi\int_{c}^{d} \bigl[b^{2} - g(y)^{2}\bigr]\,dy. \]

\begin{figureblock}
  \includegraphics[width=114.29mm]{ch08/figs/washer-split-monotone}
\figcaption{fig:8.2}{The washer reading of the shell–washer comparison, for \(f\) strictly increasing with \(0 < a\). Below \(y = c = f(a)\) every horizontal slice is the full annulus from \(a\) to \(b\); between \(c\) and \(d = f(b)\) the slice runs from \(x = g(y)\) out to \(b\). The two bands are the two terms of \(V_{\text{washer}}\).}
\end{figureblock}
Now transform the shell integral. Parts (Theorem \ref{thm:8.2}) with \(u = f(x)\) and \(dv = 2x\,dx\), so \(v = x^{2}\), gives

\[ \int_{a}^{b} 2x\,f(x)\,dx = \Bigl[x^{2}f(x)\Bigr]_{a}^{b} - \int_{a}^{b} x^{2}f'(x)\,dx = b^{2}d - a^{2}c - \int_{a}^{b} x^{2}f'(x)\,dx. \]

In the remaining integral, \(x^{2} = g(f(x))^{2}\), since \(g\) inverts \(f\). Substitution in a definite integral (Theorem 6.7), applied with the continuous outer function \(g(y)^{2}\) and with \(y = f(x)\), converts it to

\[ \int_{a}^{b} g\bigl(f(x)\bigr)^{2}\,f'(x)\,dx = \int_{c}^{d} g(y)^{2}\,dy. \]

Assembling the pieces, the shell volume equals \(\pi\bigl(b^{2}d - a^{2}c - \int_{c}^{d} g(y)^{2}\,dy\bigr)\). Expanding \(V_{\text{washer}}\) gives \(\pi\bigl(b^{2}c - a^{2}c + b^{2}(d - c) - \int_{c}^{d} g(y)^{2}\,dy\bigr)\), which is the same number, since \(b^{2}c + b^{2}(d - c) = b^{2}d\). The two methods agree, and the agreement is a computation with parts, exactly as promised. In Example 7.12’s case the common value is \(\pi\bigl(1 - 0 - \int_{0}^{1} y\,dy\bigr) = \tfrac{\pi}{2}\), the number §7.3 found twice. A strictly decreasing boundary is the mirror computation, with the annulus and the graph trading roles. For a solid whose boundary is not monotone, as in several of §7.3’s other examples, the argument above does not apply; the general statement belongs to the change of variables of Chapter 15, as §7.3 recorded.

Integration by parts is the product rule read backwards: it trades \(\int u\,dv\) for \(\int v\,du\) at the cost of the boundary product \(uv\). The trade pays when \(u\) simplifies under differentiation, and Strategy \ref{strat:8.1} makes that test explicit. Theorem \ref{thm:8.1} states the identity for antiderivatives, Theorem \ref{thm:8.2} carries it to definite integrals, and the examples display its three working patterns: a power times an integrable factor, a lone logarithm or inverse function with \(dv = dx\), and a self-reproducing integrand solved for algebraically. The same trade settled last chapter’s shell–washer question. The next section turns to integrands built entirely from trigonometric functions, where the tools are identities rather than trades, and where the reduction formula of Example \ref{ex:8.6} finds its regular use.


\sechead{8.2}{Trigonometric Integrals}
Consider \(\int \sin^{3}x\,dx\). No derivative sits beside an inner function, so no substitution presents itself, and the integrand is not a product of unrelated factors either, so parts does not simplify it. The way forward is a property specific to trigonometric functions. They convert into one another. Split off one factor of sine and rewrite the rest by the Pythagorean identity \(\sin^{2}x = 1 - \cos^{2}x\):

\[ \sin^{3}x = \bigl(1 - \cos^{2}x\bigr)\sin x. \]

Under the integral sign, the right-hand side becomes a polynomial in \(\cos x\) times \(\sin x\,dx\), and \(\sin x\,dx\) is exactly \(-du\) for \(u = \cos x\). The substitution that seemed absent has been \emph{manufactured}: one factor was spent as the differential, and an identity converted everything that remained into the new variable. That single move runs through this whole section: spend a factor and convert the rest. When no factor can be spared, a second family of identities lowers the powers instead. The identities themselves were assembled in Appendix A.2, which promised they would be needed when the techniques of integration arrived. They are needed now.

\subsechead{Powers of sine and cosine}
For an integrand \(\sin^{m}x\cos^{n}x\), the parity of the exponents (whether each is odd or even) decides the move.

\begin{envstrategy}{Strategy}{strat:8.2}{Integrating \(\sin^m x\,\cos^n x\)}
\begin{steps}
  \item \emph{If the power of sine is odd,} save one factor of \(\sin x\) for the differential and convert the remaining even power by \(\sin^{2}x = 1 - \cos^{2}x\). Then substitute \(u = \cos x\), \(du = -\sin x\,dx\).
  \item \emph{If the power of cosine is odd,} save one factor of \(\cos x\) and convert the rest by \(\cos^{2}x = 1 - \sin^{2}x\). Then substitute \(u = \sin x\), \(du = \cos x\,dx\).
  \item \emph{If both powers are even,} no factor can be spared: setting one aside would leave an odd power, which the Pythagorean identity cannot turn into a polynomial in the other function. Lower the degree with the power-reduction identities of §A.2, \(\sin^{2}x = \tfrac{1}{2}(1 - \cos 2x)\) and \(\cos^{2}x = \tfrac{1}{2}(1 + \cos 2x)\), and repeat as needed.
\end{steps}
\end{envstrategy}
\begin{workedexample}
\begin{envexample}{Example}{ex:8.7}{}
Find \(\displaystyle\int \sin^{3}x\,\cos^{2}x\,dx\).
\end{envexample}
\begin{envsolution}{Solution}{}{}
The power of sine is odd. Save one factor of \(\sin x\) and convert the remaining \(\sin^{2}x\):

\[ \int \sin^{3}x\,\cos^{2}x\,dx = \int \bigl(1 - \cos^{2}x\bigr)\cos^{2}x\,\sin x\,dx. \]

Substitute \(u = \cos x\), so that \(du = -\sin x\,dx\) and \(\sin x\,dx = -du\) (Theorem 6.6):

\[ \int \bigl(1 - u^{2}\bigr)u^{2}\,(-du) = \int \bigl(u^{4} - u^{2}\bigr)\,du = \frac{u^{5}}{5} - \frac{u^{3}}{3} + C. \]

Substituting \(u = \cos x\) back,

\[ \int \sin^{3}x\,\cos^{2}x\,dx = \frac{\cos^{5}x}{5} - \frac{\cos^{3}x}{3} + C. \]

Splitting one factor off the odd power is what made this work: that factor of sine became the differential, and everything else became a polynomial, which integrates term by term. \qedmark
\end{envsolution}
\end{workedexample}
\begin{envcaution}{Caution}{}{}
The substitution \(u = \cos x\) gives \(du = -\sin x\,dx\). The minus sign is the most commonly dropped symbol in this section. Write the line \(\sin x\,dx = -du\) explicitly before substituting, and the sign survives.
\end{envcaution}
The parity test in Strategy \ref{strat:8.2} is a test of which factor can be spent as the differential, and the integrand of Example \ref{ex:8.7} shows where the other choice leads. Spending a factor of \(\cos x\) as the differential, with \(u = \sin x\) and \(du = \cos x\,dx\), leaves \(\int u^{3}\cos x\,du\), in which the substitution is not yet complete: the old variable is still present. And a lone \(\cos x\) is not a function of \(\sin x\) until a sign is chosen: on an interval where \(\cos x \ge 0\) it equals \(\sqrt{1 - u^{2}}\), and the integral becomes \(\int u^{3}\sqrt{1 - u^{2}}\,du\), a radical instead of a polynomial. Spending a factor from the odd power leaves an even power behind, and only an even power converts by the Pythagorean identity.

\begin{workedexample}
\begin{envexample}{Example}{ex:8.8}{}
Find \(\displaystyle\int \sin^{4}x\,dx\).
\end{envexample}
\begin{envsolution}{Solution}{}{}
The power is even, so no factor can be spared and step 3 of Strategy \ref{strat:8.2} applies. Using \(\sin^{2}x = \tfrac{1}{2}(1 - \cos 2x)\) from §A.2 and squaring,

\[ \sin^{4}x = \Bigl(\frac{1 - \cos 2x}{2}\Bigr)^{2} = \frac{1}{4}\bigl(1 - 2\cos 2x + \cos^{2}2x\bigr). \]

The last term is again an even power, and the companion identity \(\cos^{2}\theta = \tfrac{1}{2}(1 + \cos 2\theta)\), applied with \(\theta = 2x\), lowers it too:

\[ \sin^{4}x = \frac{1}{4}\Bigl(1 - 2\cos 2x + \frac{1 + \cos 4x}{2}\Bigr) = \frac{3}{8} - \frac{\cos 2x}{2} + \frac{\cos 4x}{8}. \]

Every term is now a plain cosine, and integration is immediate:

\[ \int \sin^{4}x\,dx = \frac{3x}{8} - \frac{\sin 2x}{4} + \frac{\sin 4x}{32} + C. \]

The reduction formula of Example \ref{ex:8.6} reaches the same answer by a different route. It expresses \(\int\sin^{4}x\,dx\) in terms of \(\int\sin^{2}x\,dx\), and the two results agree after the double-angle identity is applied. When several powers must be integrated in sequence, the reduction formula is the more systematic tool; for a single even power, power reduction is the shorter one. \qedmark
\end{envsolution}
\end{workedexample}
\begin{figureblock}
  \includegraphics[width=122.22mm]{ch08/figs/sin-squared-midline}
\figcaption{fig:8.3}{The graph of \(\sin^{2}x\) oscillates about the midline \(y = \tfrac{1}{2}\) with period \(\pi\): the identity’s constant term is the average height, and each valley exactly fills a peak. The zeros touch the axis without crossing it.}
\end{figureblock}
\subsechead{Tangents and secants}
Recall from Definition A.1 that \(\sec x = \tfrac{1}{\cos x}\), \(\csc x = \tfrac{1}{\sin x}\), and \(\cot x = \tfrac{\cos x}{\sin x}\). For integration, tangent and secant form a closed pair. Their derivatives stay in the pair, since \(\tfrac{d}{dx}\tan x = \sec^{2}x\) and \(\tfrac{d}{dx}\sec x = \sec x\tan x\) (Example 3.2), and the Pythagorean companion \(1 + \tan^{2}x = \sec^{2}x\) (Proposition A.1) converts between them. So the same manufacturing move applies. Spend \(\sec^{2}x\,dx\) as the differential of \(\tan x\), or spend \(\sec x\tan x\,dx\) as the differential of \(\sec x\), and convert what remains.

\begin{envstrategy}{Strategy}{strat:8.3}{Integrating \(\tan^m x\,\sec^n x\)}
\begin{steps}
  \item \emph{If the power of secant is even} (\(n \ge 2\)), save a factor of \(\sec^{2}x\) and convert the remaining secants by \(\sec^{2}x = 1 + \tan^{2}x\). Then substitute \(u = \tan x\), \(du = \sec^{2}x\,dx\).
  \item \emph{If the power of tangent is odd} (and \(n \ge 1\)), save a factor of \(\sec x\tan x\) and convert the remaining tangents by \(\tan^{2}x = \sec^{2}x - 1\). Then substitute \(u = \sec x\), \(du = \sec x\tan x\,dx\).
  \item \emph{If neither case applies,} convert tangents by \(\tan^{2}x = \sec^{2}x - 1\) and fall back on parts and the basic antiderivatives of Proposition \ref{prop:8.1} below, as in Example \ref{ex:8.10}. A lone even power of tangent splits this way into lower powers plus a constant; a lone odd power reduces the same way to \(\int\tan x\,dx\) itself, which Proposition \ref{prop:8.1} supplies.
\end{steps}
\end{envstrategy}
\begin{workedexample}
\begin{envexample}{Example}{ex:8.9}{}
Find \(\displaystyle\int \tan^{6}x\,\sec^{4}x\,dx\).
\end{envexample}
\begin{envsolution}{Solution}{}{}
The power of secant is even. Save one factor of \(\sec^{2}x\) and convert the other by \(\sec^{2}x = 1 + \tan^{2}x\):

\[ \int \tan^{6}x\,\sec^{4}x\,dx = \int \tan^{6}x\,\bigl(1 + \tan^{2}x\bigr)\,\sec^{2}x\,dx. \]

Substitute \(u = \tan x\), \(du = \sec^{2}x\,dx\):

\[ \int u^{6}\bigl(1 + u^{2}\bigr)\,du = \int \bigl(u^{6} + u^{8}\bigr)\,du = \frac{u^{7}}{7} + \frac{u^{9}}{9} + C. \]

Substituting back,

\[ \int \tan^{6}x\,\sec^{4}x\,dx = \frac{\tan^{7}x}{7} + \frac{\tan^{9}x}{9} + C, \]

valid on any interval where \(\tan x\) and \(\sec x\) are defined. \qedmark
\end{envsolution}
\end{workedexample}
Strategy \ref{strat:8.3} stands on two basic antiderivatives that its case analysis cannot produce. Its third step reduces leftover tangents to \(\int\tan x\,dx\) itself, and the route of Example \ref{ex:8.10} runs through \(\int\sec x\,dx\). The antiderivative table of §6.4 has neither: its trigonometric rows are the Chapter 3 derivatives read backwards, so \(\sec^{2}x\) appears there but \(\tan x\) itself does not. The next proposition supplies both, together with their two companions.

\begin{envproposition}{Proposition}{prop:8.1}{Antiderivatives of the remaining trigonometric functions}
On any interval where the integrand is continuous,

\[ \begin{aligned}
        \int \tan x\,dx &= \ln\lvert\sec x\rvert + C, \\[2pt]
        \int \cot x\,dx &= \ln\lvert\sin x\rvert + C, \\[2pt]
        \int \sec x\,dx &= \ln\lvert\sec x + \tan x\rvert + C, \\[2pt]
        \int \csc x\,dx &= -\ln\lvert\csc x + \cot x\rvert + C.
      \end{aligned} \]
\end{envproposition}
\begin{envproof}{Proof}{}{}
For the tangent, write \(\tan x = \tfrac{\sin x}{\cos x}\) and substitute \(u = \cos x\), \(du = -\sin x\,dx\) (Theorem 6.6). Using the antiderivative \(\int\tfrac{du}{u} = \ln\lvert u\rvert + C\) from §6.4,

\[ \int \tan x\,dx = \int \frac{\sin x}{\cos x}\,dx = -\int \frac{du}{u} = -\ln\lvert\cos x\rvert + C = \ln\lvert\sec x\rvert + C, \]

where the last step uses \(-\ln t = \ln\tfrac{1}{t}\) for \(t > 0\). The cotangent line is the same computation with \(u = \sin x\), this time without the minus sign.

For the secant no substitution presents itself, and the standard route is a device that was found by trial, not derived: multiply and divide by \(\sec x + \tan x\). Then

\[ \int \sec x\,dx = \int \frac{\sec^{2}x + \sec x\tan x}{\sec x + \tan x}\,dx, \]

and the numerator is exactly the derivative of the denominator, by Example 3.2. With \(u = \sec x + \tan x\) the integral is \(\int\tfrac{du}{u} = \ln\lvert u\rvert + C\), which is the stated formula. Whatever one thinks of the trick, the result is verified honestly by differentiating:

\[ \begin{aligned}
        \frac{d}{dx}\,\ln\lvert\sec x + \tan x\rvert &= \frac{\sec x\tan x + \sec^{2}x}{\sec x + \tan x} \\[2pt]
          &= \frac{\sec x\,\bigl(\tan x + \sec x\bigr)}{\sec x + \tan x} = \sec x.
      \end{aligned} \]

The cosecant line follows by the same verification, differentiating \(-\ln\lvert\csc x + \cot x\rvert\) and simplifying with \(\tfrac{d}{dx}\csc x = -\csc x\cot x\) and \(\tfrac{d}{dx}\cot x = -\csc^{2}x\). \qedmark
\end{envproof}
Each formula holds on every interval avoiding the integrand’s singularities, and the constant \(C\) may differ from interval to interval, as with \(\ln\lvert x\rvert\) in §6.4.

The antiderivative of the secant was wanted before there was a calculus to supply it. On the Mercator chart of 1569, the north-south scale has to stretch with latitude in proportion to the secant of the latitude, so the map distance out to latitude \(\varphi\) is \(\int_{0}^{\varphi}\sec\theta\,d\theta\). Edward Wright published a table of those values in 1599, obtained by adding secants in small steps, since no formula for them was known. Around 1645 Henry Bond noticed that Wright’s table agreed with a table of logarithms of tangents and conjectured a closed form on that evidence alone; James Gregory proved the conjecture in 1668, and Isaac Barrow gave a clearer proof in 1670. The device above looks unmotivated because the order of events was the reverse of the usual one: the formula was found first and justified afterwards.

\begin{workedexample}
\begin{envexample}{Example}{ex:8.10}{}
Find \(\displaystyle\int \sec^{3}x\,dx\).
\end{envexample}
\begin{envsolution}{Solution}{}{}
The power of secant is odd and no tangent is present, so neither case of Strategy \ref{strat:8.3} applies. Following its third step, we use parts (Theorem \ref{thm:8.1}). Take \(u = \sec x\) and \(dv = \sec^{2}x\,dx\), so that \(du = \sec x\tan x\,dx\) by Example 3.2 and \(v = \tan x\):

\[ \int \sec^{3}x\,dx = \sec x\tan x - \int \sec x\tan^{2}x\,dx. \]

Convert the tangents by \(\tan^{2}x = \sec^{2}x - 1\) (Proposition A.1):

\[ \int \sec^{3}x\,dx = \sec x\tan x - \int \sec^{3}x\,dx + \int \sec x\,dx. \]

The original integral has reappeared, and we solve for it as in Example \ref{ex:8.4}. Adding \(\int\sec^{3}x\,dx\) to both sides, dividing by \(2\), and evaluating \(\int\sec x\,dx\) by Proposition \ref{prop:8.1}:

\[ \int \sec^{3}x\,dx = \frac{1}{2}\Bigl(\sec x\tan x + \ln\lvert\sec x + \tan x\rvert\Bigr) + C. \]

This integral appears throughout the next section, where substitutions of the form \(x = a\tan\theta\) produce it repeatedly. The mirror integral \(\int\csc^{3}x\,dx\) is computed by the same two moves, parts and then Proposition \ref{prop:8.1}’s cosecant formula. \qedmark
\end{envsolution}
\end{workedexample}
The cotangent-cosecant family behaves in exactly the same way, with the same split into cases. For \(\int\cot^{m}x\,\csc^{n}x\,dx\): if \(n\) is even, save a factor \(\csc^{2}x\,dx\), which is the differential of \(-\cot x\), and convert what remains by \(1 + \cot^{2}x = \csc^{2}x\) (Proposition A.1); if \(m\) is odd, save a factor \(\csc x\cot x\,dx\), which is the differential of \(-\csc x\), and convert what remains by the same identity. The only change from Strategy \ref{strat:8.3} is a minus sign carried through each substitution.

\subsechead{Products with different frequencies}
The last family is \(\int\sin mx\cos nx\,dx\) and its siblings, where the two factors oscillate at different rates. No substitution connects them. Here the product-to-sum identities do all the work in one line, by replacing the product with a sum of single trigonometric functions. Proposition A.2’s first line is the one needed for a sine times a cosine:

\[ \sin u\cos v = \tfrac{1}{2}\bigl[\sin(u + v) + \sin(u - v)\bigr]. \]

\begin{workedexample}
\begin{envexample}{Example}{ex:8.11}{}
Find \(\displaystyle\int \sin 4x\,\cos 5x\,dx\).
\end{envexample}
\begin{envsolution}{Solution}{}{}
Apply the identity with \(u = 4x\) and \(v = 5x\):

\[ \sin 4x\,\cos 5x = \tfrac{1}{2}\bigl[\sin 9x + \sin(-x)\bigr] = \tfrac{1}{2}\bigl[\sin 9x - \sin x\bigr], \]

since sine is an odd function, \(\sin(-x) = -\sin x\). The product has become a sum of two sines at fixed frequencies, and each integrates immediately:

\[ \int \sin 4x\,\cos 5x\,dx = \frac{1}{2}\Bigl(-\frac{\cos 9x}{9} + \cos x\Bigr) + C = \frac{\cos x}{2} - \frac{\cos 9x}{18} + C. \]

For a product of two sines, or of two cosines, the other two lines of Proposition A.2 convert the product the same way. \qedmark
\end{envsolution}
\end{workedexample}
This section added no new integration principle. Its engine is still substitution, fed by identities: one factor of an odd power is spent as the differential, the Pythagorean identities convert what remains, power reduction lowers even powers, and product-to-sum splits mixed frequencies. The four antiderivatives of Proposition \ref{prop:8.1} fill the last gaps in the basic trigonometric table, and Example \ref{ex:8.10}’s \(\int\sec^{3}x\,dx\) is recorded for immediate reuse. The next section turns this machinery outward, onto integrands that contain no trigonometry at all: radicals such as \(\sqrt{9 - x^{2}}\), which a well-chosen trigonometric substitution converts into exactly the integrals this section taught.


\sechead{8.3}{Trigonometric Substitution}
The radical \(\sqrt{9 - x^{2}}\) appeared in §8.2’s closing sentence as the next target, and it resists every method so far. No substitution of the form \(u = 9 - x^{2}\) helps unless a spare factor of \(x\) happens to sit in the integrand to supply \(du\). When it does not, the root over a sum or difference of squares blocks every algebraic move. What removes it is a change of viewpoint. The Pythagorean identity is itself a statement about sums of squares: setting \(x = 3\sin\theta\) turns \(9 - x^{2}\) into \(9\cos^{2}\theta\), a perfect square, and the root falls away. The substitution changes the variable, so a return step waits at the end; what is gained is an integrand made of trigonometric functions, which §8.2 taught us to handle.

One point of honesty before we state the method. In Chapter 6 we substituted \(u = g(x)\), replacing the variable by an expression. Here we run the Substitution Rule (Theorem 6.6) in the other direction, replacing the variable \(x\) by a formula \(x = g(\theta)\) in a new variable. For this to be reversible, \(g\) must be one-to-one. That is what the restricted ranges in the table below arrange: on \([-\tfrac{\pi}{2}, \tfrac{\pi}{2}]\) the function \(a\sin\theta\) takes each value exactly once, so \(\theta = \arcsin\tfrac{x}{a}\) is well defined, just as in §1.2. The same restriction settles a second question. A square root such as \(\sqrt{9\cos^{2}\theta}\) equals \(3\lvert\cos\theta\rvert\), and on \([-\tfrac{\pi}{2}, \tfrac{\pi}{2}]\) the cosine is nonnegative, so the absolute value can be dropped with justification rather than by hope. Every solution in this section begins by naming its substitution, its \(\theta\)-range, and the sign that the range fixes.

Throughout, \(a\) denotes a positive constant: the radicals involve only \(a^{2}\), so nothing is lost by taking \(a\) positive. Three radical patterns cover the standard cases, one for each rearrangement of the Pythagorean identity (Proposition A.1).

\begin{envstrategy}{Strategy}{strat:8.4}{Trigonometric substitution}
First check for an ordinary substitution: if the integrand carries a spare factor of \(x\), as in \(\int x\sqrt{1 - x^{2}}\,dx\), then \(u = 1 - x^{2}\) finishes without trigonometry. The same check returns as step 2 of §8.5’s four-step strategy. Otherwise, match the radical to a row:

\begin{steps}
  \item \(\sqrt{a^{2} - x^{2}}\): substitute \(x = a\sin\theta\), \(\theta \in \bigl[-\tfrac{\pi}{2}, \tfrac{\pi}{2}\bigr]\). Then \(1 - \sin^{2}\theta = \cos^{2}\theta\) gives \(\sqrt{a^{2} - x^{2}} = a\cos\theta\), with \(\cos\theta \ge 0\) on the range.
  \item \(\sqrt{a^{2} + x^{2}}\): substitute \(x = a\tan\theta\), \(\theta \in \bigl(-\tfrac{\pi}{2}, \tfrac{\pi}{2}\bigr)\). Then \(1 + \tan^{2}\theta = \sec^{2}\theta\) gives \(\sqrt{a^{2} + x^{2}} = a\sec\theta\), with \(\sec\theta > 0\) on the range.
  \item \(\sqrt{x^{2} - a^{2}}\): substitute \(x = a\sec\theta\), \(\theta \in \bigl[0, \tfrac{\pi}{2}\bigr)\) for the component \(x \ge a\) of the domain. Then \(\sec^{2}\theta - 1 = \tan^{2}\theta\) gives \(\sqrt{x^{2} - a^{2}} = a\tan\theta\), with \(\tan\theta \ge 0\) on the range. (For the component \(x \le -a\), see the caution at the end of the section.)
\end{steps}
\begin{informal}
To return to \(x\), draw a right triangle that matches the substitution: mark the angle \(\theta\), and label the sides so that the substitution’s ratio holds. The remaining side, by the Pythagorean theorem, is the radical. Every trigonometric function of \(\theta\) can then be read off as a ratio of labeled sides. The triangle is a mnemonic for identities that the \(\theta\)-range makes valid.
\end{informal}
\end{envstrategy}
\begin{figureblock}
  \includegraphics[width=47.09mm]{ch08/figs/trig-sub-triangles-1}\hspace{3.86mm}\includegraphics[width=47.09mm]{ch08/figs/trig-sub-triangles-2}\hspace{3.86mm}\includegraphics[width=47.09mm]{ch08/figs/trig-sub-triangles-3}
\figcaption{fig:8.4}{The three reference triangles of Strategy \ref{strat:8.4}, one per substitution. Label the sides so that the substitution’s ratio holds; the remaining side, by the Pythagorean theorem, is the radical, and every trigonometric function of \(\theta\) can be read off as a ratio of sides. The triangles are a mnemonic for identities that the stated \(\theta\)-ranges make valid.}
\end{figureblock}
\begin{workedexample}
\begin{envexample}{Example}{ex:8.12}{}
Find \(\displaystyle\int \frac{\sqrt{9 - x^{2}}}{x^{2}}\,dx\) for \(0 < x < 3\).
\end{envexample}
\begin{envsolution}{Solution}{}{}
The radical matches row 1 with \(a = 3\). Substitute \(x = 3\sin\theta\); on the component \(0 < x < 3\) the range narrows to \(\theta \in \bigl(0, \tfrac{\pi}{2}\bigr)\), where \(\cos\theta > 0\), so \(\sqrt{9 - x^{2}} = 3\cos\theta\) and \(dx = 3\cos\theta\,d\theta\). The integral becomes

\[ \int \frac{3\cos\theta}{9\sin^{2}\theta}\,3\cos\theta\,d\theta = \int \frac{\cos^{2}\theta}{\sin^{2}\theta}\,d\theta = \int \cot^{2}\theta\,d\theta. \]

By the Pythagorean companion \(1 + \cot^{2}\theta = \csc^{2}\theta\) (Proposition A.1), the integrand is \(\csc^{2}\theta - 1\), and \(\int\csc^{2}\theta\,d\theta = -\cot\theta + C\) since \(\tfrac{d}{d\theta}\cot\theta = -\csc^{2}\theta\). So

\[ \int \cot^{2}\theta\,d\theta = -\cot\theta - \theta + C. \]

Now return to \(x\). In the right triangle for \(\sin\theta = \tfrac{x}{3}\), the hypotenuse is \(3\), the side opposite \(\theta\) is \(x\), and the remaining side is \(\sqrt{9 - x^{2}}\). Reading off the ratio, \(\cot\theta = \tfrac{\sqrt{9 - x^{2}}}{x}\), and \(\theta = \arcsin\tfrac{x}{3}\) directly. Therefore

\[ \int \frac{\sqrt{9 - x^{2}}}{x^{2}}\,dx = -\frac{\sqrt{9 - x^{2}}}{x} - \arcsin\frac{x}{3} + C \qquad\text{on } (0, 3). \]

The same computation serves on \((-3, 0)\), where \(\theta \in \bigl(-\tfrac{\pi}{2}, 0\bigr)\) and the cosine is still positive; only the component changes. \qedmark
\end{envsolution}
\end{workedexample}
\begin{workedexample}
\begin{envexample}{Example}{ex:8.13}{}
Find the area enclosed by the ellipse \(\dfrac{x^{2}}{a^{2}} + \dfrac{y^{2}}{b^{2}} = 1\), where \(a, b > 0\).
\end{envexample}
\begin{envsolution}{Solution}{}{}
Solving for the upper half gives \(y = \tfrac{b}{a}\sqrt{a^{2} - x^{2}}\) for \(-a \le x \le a\), so the enclosed area is twice the integral of that height:

\[ A = 2\,\frac{b}{a}\int_{-a}^{a} \sqrt{a^{2} - x^{2}}\,dx = 4\,\frac{b}{a}\int_{0}^{a} \sqrt{a^{2} - x^{2}}\,dx, \]

where the second equality uses the evenness of the integrand (Theorem 6.8). The radical matches row 1. Substitute \(x = a\sin\theta\), \(dx = a\cos\theta\,d\theta\), and transform the limits by Theorem 6.7: \(x = 0\) gives \(\theta = 0\) and \(x = a\) gives \(\theta = \tfrac{\pi}{2}\). On \(\bigl[0, \tfrac{\pi}{2}\bigr]\), where \(\cos\theta \ge 0\), \(\sqrt{a^{2} - x^{2}} = a\cos\theta\), so

\[ \begin{aligned}
          \int_{0}^{a} \sqrt{a^{2} - x^{2}}\,dx &= \int_{0}^{\pi/2} a^{2}\cos^{2}\theta\,d\theta = a^{2}\int_{0}^{\pi/2} \frac{1 + \cos 2\theta}{2}\,d\theta \\[2pt]
            &= a^{2}\left[\frac{\theta}{2} + \frac{\sin 2\theta}{4}\right]_{0}^{\pi/2} = \frac{\pi a^{2}}{4},
        \end{aligned} \]

using the power-reduction identity \(\cos^{2}\theta = \tfrac{1}{2}(1 + \cos 2\theta)\) of §A.2. Therefore

\[ A = 4\,\frac{b}{a}\cdot\frac{\pi a^{2}}{4} = \pi a b. \]

For \(a = b = r\) the ellipse is a circle and the formula returns \(\pi r^{2}\), as it must. \qedmark
\end{envsolution}
\end{workedexample}
The ellipse’s area is now settled; its boundary length is not. For \(a \ne b\), the arc length integral of §7.6 for this curve contains a radical that no row of Strategy \ref{strat:8.4}, and no other technique of this chapter, can turn into an integrand whose antiderivative is built from the functions we have named. (The circle \(a = b = r\) is the exception: there the integrand reduces to a constant times \(1/\sqrt{r^{2} - x^{2}}\), which row 1 converts in one line.) The area took four lines; the perimeter sits outside this chapter’s reach. Section 8.7 returns to exactly this kind of integral with numerical rules.

\begin{figureblock}
  \includegraphics[width=101.06mm]{ch08/figs/ellipse-quarter-area}
\figcaption{fig:8.5}{The quarter-ellipse region behind Example \ref{ex:8.13}. Solving for the upper half doubles the region once, evenness doubles it again, so the enclosed area is four times the shaded first-quadrant integral. The intercepts sit at \((a, 0)\) and \((0, b)\).}
\end{figureblock}
\begin{workedexample}
\begin{envexample}{Example}{ex:8.14}{}
Find \(\displaystyle\int \frac{dx}{x^{2}\sqrt{x^{2} + 4}}\).
\end{envexample}
\begin{envsolution}{Solution}{}{}
The radical matches row 2 with \(a = 2\). Substitute \(x = 2\tan\theta\), \(\theta \in \bigl(-\tfrac{\pi}{2}, \tfrac{\pi}{2}\bigr)\); we work on the component \(x > 0\), where \(\theta \in \bigl(0, \tfrac{\pi}{2}\bigr)\). Then \(dx = 2\sec^{2}\theta\,d\theta\), and since \(\sec\theta > 0\) on the range, \(\sqrt{x^{2} + 4} = 2\sec\theta\). The integral becomes

\[ \int \frac{2\sec^{2}\theta}{4\tan^{2}\theta\cdot 2\sec\theta}\,d\theta = \frac{1}{4}\int \frac{\sec\theta}{\tan^{2}\theta}\,d\theta = \frac{1}{4}\int \frac{\cos\theta}{\sin^{2}\theta}\,d\theta, \]

after writing everything in sines and cosines. With \(u = \sin\theta\), \(du = \cos\theta\,d\theta\), this is \(\tfrac{1}{4}\int u^{-2}\,du = -\tfrac{1}{4u} + C = -\tfrac{1}{4}\csc\theta + C\). In the triangle for \(\tan\theta = \tfrac{x}{2}\), the side opposite \(\theta\) is \(x\), the adjacent side is \(2\), and the hypotenuse is \(\sqrt{x^{2} + 4}\), so \(\csc\theta = \tfrac{\sqrt{x^{2} + 4}}{x}\). Therefore

\[ \int \frac{dx}{x^{2}\sqrt{x^{2} + 4}} = -\frac{\sqrt{x^{2} + 4}}{4x} + C. \]

Differentiating the answer shows it is an antiderivative on \(x < 0\) as well, so the one formula serves both components. \qedmark
\end{envsolution}
\end{workedexample}
When the radical stands alone in the integrand, as in \(\int\sqrt{x^{2} + a^{2}}\,dx\), the same tangent substitution leads to \(\int\sec^{3}\theta\,d\theta\), which Example \ref{ex:8.10} supplies; this book leaves that longer computation unworked.

\begin{workedexample}
\begin{envexample}{Example}{ex:8.15}{}
Find \(\displaystyle\int \frac{dx}{\sqrt{x^{2} - a^{2}}}\) for \(x > a > 0\).
\end{envexample}
\begin{envsolution}{Solution}{}{}
The radical matches row 3. Substitute \(x = a\sec\theta\), \(\theta \in \bigl(0, \tfrac{\pi}{2}\bigr)\) on this component; then \(dx = a\sec\theta\tan\theta\,d\theta\), and since \(\tan\theta > 0\) on the range, \(\sqrt{x^{2} - a^{2}} = a\tan\theta\). The integral collapses:

\[ \int \frac{a\sec\theta\tan\theta}{a\tan\theta}\,d\theta = \int \sec\theta\,d\theta = \ln\lvert\sec\theta + \tan\theta\rvert + C, \]

by Proposition \ref{prop:8.1}. In the triangle for \(\sec\theta = \tfrac{x}{a}\), the hypotenuse is \(x\), the adjacent side is \(a\), and the opposite side is \(\sqrt{x^{2} - a^{2}}\), so

\[ \int \frac{dx}{\sqrt{x^{2} - a^{2}}} = \ln\left\lvert \frac{x}{a} + \frac{\sqrt{x^{2} - a^{2}}}{a} \right\rvert + C = \ln\bigl(x + \sqrt{x^{2} - a^{2}}\bigr) + C_{1}, \]

where \(\ln a\) was absorbed into the constant, and the absolute value was dropped because \(x > a > 0\) makes the argument positive on this component. \qedmark
\end{envsolution}
\end{workedexample}
\begin{envcaution}{Caution}{}{}
The domain of \(\sqrt{x^{2} - a^{2}}\) has two components, and row 3’s range covers only \(x \ge a\). On the other component, \(x \le -a\), use the alternate range \(\theta \in \bigl(\tfrac{\pi}{2}, \pi\bigr]\): there \(\tan\theta \le 0\), so the resolved radical is \(\sqrt{x^{2} - a^{2}} = -a\tan\theta\), and the extra sign travels through the computation. The shortcut is to check the final formula directly. Differentiating \(\ln\bigl\lvert x + \sqrt{x^{2} - a^{2}}\bigr\rvert\) gives \(\tfrac{1}{\sqrt{x^{2} - a^{2}}}\) on \emph{both} components, so Example \ref{ex:8.15}’s answer, with the absolute value restored, serves everywhere. An antiderivative claimed across a two-component domain must be checked on each component.
\end{envcaution}
\begin{workedexample}
\begin{envexample}{Example}{ex:8.16}{}
Find \(\displaystyle\int \frac{x}{\sqrt{3 - 2x - x^{2}}}\,dx\).
\end{envexample}
\begin{envsolution}{Solution}{}{}
No row of Strategy \ref{strat:8.4} matches until the quadratic under the root is recentred. Completing the square, as in the closing paragraph of §A.5, with \(x^{2} + 2x = (x + 1)^{2} - 1\):

\[ 3 - 2x - x^{2} = 4 - (x + 1)^{2}, \]

which is positive exactly for \(-3 < x < 1\), the domain of the integrand. The shift \(u = x + 1\), \(du = dx\), \(x = u - 1\) turns the integral into row-1 form:

\[ \int \frac{u - 1}{\sqrt{4 - u^{2}}}\,du = \int \frac{u}{\sqrt{4 - u^{2}}}\,du - \int \frac{du}{\sqrt{4 - u^{2}}}. \]

The first piece has the spare factor that Strategy \ref{strat:8.4}’s lead-in asks about: \(w = 4 - u^{2}\), \(dw = -2u\,du\) gives \(-\tfrac{1}{2}\int w^{-1/2}\,dw = -\sqrt{4 - u^{2}}\). The second piece is row 1 with \(a = 2\) and nothing else in the integrand: with \(u = 2\sin\varphi\), \(\varphi \in \bigl(-\tfrac{\pi}{2}, \tfrac{\pi}{2}\bigr)\), where \(\cos\varphi > 0\), both \(du = 2\cos\varphi\,d\varphi\) and \(\sqrt{4 - u^{2}} = 2\cos\varphi\) reduce it to \(\int d\varphi = \varphi = \arcsin\tfrac{u}{2}\). Assembling and returning to \(x\):

\[ \int \frac{x}{\sqrt{3 - 2x - x^{2}}}\,dx = -\sqrt{3 - 2x - x^{2}} - \arcsin\frac{x + 1}{2} + C, \]

valid on \((-3, 1)\). Completing the square is what §A.5 promised it would be here: the one algebraic step that brings a general quadratic radical within reach of the table. \qedmark
\end{envsolution}
\end{workedexample}
Completing the square clears the linear term; a leading coefficient other than \(1\) needs one more preliminary step. The three rows of Strategy \ref{strat:8.4} are written with \(x^{2}\) by itself, so the leading coefficient must be factored out first. Then \(\sqrt{9x^{2} - 4} = 3\sqrt{x^{2} - \bigl(\tfrac{2}{3}\bigr)^{2}}\) is row 3 with \(a = \tfrac{2}{3}\), and \(\sqrt{9 - 4x^{2}} = 2\sqrt{\bigl(\tfrac{3}{2}\bigr)^{2} - x^{2}}\) is row 1 with \(a = \tfrac{3}{2}\). A linear substitution does the same work, as in Example \ref{ex:8.21} of the next section, where \(u = 2x - 1\) clears a leading coefficient from a quadratic denominator.

\begin{envcaution}{Caution}{}{}
For a definite integral, either transform the limits along with the variable (Theorem 6.7), as in Example \ref{ex:8.13}, or back-substitute to \(x\) first and then apply the original limits. Do not mix the two: evaluating a \(\theta\)-antiderivative at \(x\)-limits is the classic error of this method.
\end{envcaution}
Trigonometric substitution replaces the variable by a trigonometric formula so that a Pythagorean identity turns the quadratic under the root into a perfect square. The discipline is in the ranges: each substitution is one-to-one on its stated interval, the interval fixes the radical’s sign, and the labeled triangle carries the answer back to \(x\). With completing the square, the method reaches any \(\sqrt{px^{2} + qx + r}\). What it does not touch is the other standard shape built from polynomials: a ratio of two of them, with no radical in sight. Integrating rational functions is the next section’s business, and the algebra for it is already waiting in Appendix A.4.


\sechead{8.4}{Integration of Rational Functions by Partial Fractions}
A rational function, a ratio of polynomials, matches nothing in the chapter so far. The integrand \(\dfrac{2x^{2} - x + 4}{x^{3} + 4x}\) offers no product to trade by parts, no trigonometry, and no radical. Yet one of the smallest rational integrands is already easy, because §A.4 opened with the identity

\[ \frac{1}{x(x + 1)} = \frac{1}{x} - \frac{1}{x + 1}, \]

and each piece on the right integrates on sight, to \(\ln\lvert x\rvert - \ln\lvert x + 1\rvert + C\) by §6.4. Appendix A.4 built the whole algebra of such splittings and remarked that its real payoff would come in integration. This section collects that payoff. Every proper rational function is a sum of fragments of a few standard shapes (Proposition A.7), the constants are found by the routine of Strategy A.2, and each fragment shape has a standard antiderivative. For rational integrands, unlike most others, the method always succeeds in producing the decomposition, and every fragment it produces—except one, flagged where it arises—is then integrated in full.

The claim “every proper rational function” rests on one algebraic fact that this book states but does not prove.

\begin{envcaution}{Caution}{}{}
Proposition A.7 begins by factoring the denominator completely into linear and irreducible quadratic factors. That such a factorization always exists is a consequence of the Fundamental Theorem of Algebra, whose proof lies outside this book. We take the factorization fact on credit; it is the first of two results this chapter borrows without proof, and apart from the one fragment integral left as stated below (repeated irreducible quadratic factors), everything else in the section is proved or computed in full.
\end{envcaution}
Existence is not the same as a method for finding one. The credit above, and the promise that the method always succeeds, guarantee that a factorization exists; neither says how to produce it. The first tool is the factor theorem: \(r\) is a root of a polynomial exactly when \(x - r\) divides it. When the coefficients are integers this narrows the search, because an integer root must divide the constant term, so the divisors of that term form a short list of candidates to test. Every denominator in this section factors by hand, sometimes by that test and sometimes by a shortcut: the quadratic formula, which supplies the roots of the quadratic factor in Example \ref{ex:8.18} below, or grouping, as in Example \ref{ex:8.19}. When no factorization can be found, the method stops at step 2 of Strategy \ref{strat:8.5}, and no integration technique can supply the missing factorization. From here on we take the factoring as done.

Before the fragments can be integrated, one small antiderivative is needed: the scaled version of §6.4’s arctangent row. For \(a > 0\), substitute \(u = x/a\), \(du = dx/a\), in the table entry \(\int\frac{du}{1 + u^{2}} = \arctan u + C\):

\[ \int \frac{dx}{x^{2} + a^{2}} = \frac{1}{a^{2}}\int \frac{a\,du}{1 + u^{2}} = \frac{1}{a}\arctan\frac{x}{a} + C. \tag{8.A} \]

Once the denominator is factored, Proposition A.7 determines which fragments it produces. That rule is short enough to state here. Each factor of the denominator contributes one fragment for each power up to the exponent it carries. The numerator’s degree depends on the factor itself, not on the power it is raised to: a constant over a linear factor, a linear \(Bx + C\) over an irreducible quadratic. A distinct linear factor \(x - r\) contributes \(\frac{A}{x - r}\), and a linear factor written in another form, such as \(2x - 1\), is treated the same way. A repeated linear factor \((x - r)^{k}\) contributes \(k\) fragments,

\[ \frac{A_{1}}{x - r} + \frac{A_{2}}{(x - r)^{2}} + \cdots + \frac{A_{k}}{(x - r)^{k}}, \]

one for each power. An irreducible quadratic \(x^{2} + px + q\) contributes \(\frac{Bx + C}{x^{2} + px + q}\). A repeated irreducible quadratic \((x^{2} + px + q)^{k}\) contributes one such fragment over each power up to \(k\). This section writes that template down but only describes its integration, without a worked example. What the template leaves open are the constants; finding them is step 3 of the strategy below, and Examples \ref{ex:8.18} through \ref{ex:8.20} carry that step out in full.

With (8.A) in hand, every fragment shape that Proposition A.7 can produce has a known antiderivative. A constant over a linear factor gives a logarithm, \(\int\frac{A}{x - r}\,dx = A\ln\lvert x - r\rvert + C\). A constant over a power of a linear factor integrates by the power rule, \(\int\frac{A}{(x - r)^{k}}\,dx = \frac{-A}{(k - 1)(x - r)^{k - 1}} + C\) for \(k \ge 2\). A linear numerator over an irreducible quadratic splits in two: a multiple of the quadratic’s derivative, giving a logarithm of the quadratic, plus a constant piece that (8.A) handles once the square is completed. The examples run each shape in turn; the strategy box states the pipeline.

\begin{envstrategy}{Strategy}{strat:8.5}{Integrating a rational function}
\begin{steps}
  \item \emph{If the fraction is improper, divide first.} Numerator degree at or above denominator degree means polynomial division, leaving a polynomial plus a proper remainder fraction (Strategy A.2, step 1). The polynomial integrates term by term.
  \item \emph{Factor the denominator completely} into linear and irreducible quadratic factors. A quadratic \(x^{2} + px + q\) is irreducible exactly when \(p^{2} < 4q\).
  \item \emph{Decompose.} Write the template of Proposition A.7 and find the constants by Strategy A.2. The cover-up shortcut reaches the constant over each distinct linear factor at once.
  \item \emph{Integrate fragment by fragment:} logarithms from linear factors, power-rule antiderivatives from their powers, and a logarithm plus an arctangent, via completing the square (§A.5) and (8.A), from each irreducible quadratic.
\end{steps}
\end{envstrategy}
\begin{workedexample}
\begin{envexample}{Example}{ex:8.17}{}
Find \(\displaystyle\int \frac{x^{3} + x}{x - 1}\,dx\).
\end{envexample}
\begin{envsolution}{Solution}{}{}
The numerator’s degree exceeds the denominator’s, so step 1 applies before anything else. Polynomial division gives

\[ \frac{x^{3} + x}{x - 1} = x^{2} + x + 2 + \frac{2}{x - 1}, \]

which can be checked by recombining over the common denominator. Now every piece is standard:

\[ \int \frac{x^{3} + x}{x - 1}\,dx = \frac{x^{3}}{3} + \frac{x^{2}}{2} + 2x + 2\ln\lvert x - 1\rvert + C, \]

on any interval avoiding \(x = 1\). The division step is not optional. No sum of terms \(\frac{A}{x-1}\) alone could reproduce a numerator that outgrows its denominator, which is the point of §A.4’s closing caution. \qedmark
\end{envsolution}
\end{workedexample}
\begin{workedexample}
\begin{envexample}{Example}{ex:8.18}{}
Find \(\displaystyle\int \frac{x^{2} + 2x - 1}{2x^{3} + 3x^{2} - 2x}\,dx\).
\end{envexample}
\begin{envsolution}{Solution}{}{}
The fraction is proper, so factoring comes first: \(2x^{3} + 3x^{2} - 2x = x(2x^{2} + 3x - 2) = x(2x - 1)(x + 2)\), three distinct linear factors. Proposition A.7’s template reads

\[ \frac{x^{2} + 2x - 1}{x(2x - 1)(x + 2)} = \frac{A}{x} + \frac{B}{2x - 1} + \frac{C}{x + 2}. \]

Clearing denominators,

\[ x^{2} + 2x - 1 = A(2x - 1)(x + 2) + Bx(x + 2) + Cx(2x - 1), \]

and the cover-up shortcut of Strategy A.2 evaluates at each factor’s root. Setting \(x = 0\) gives \(-1 = -2A\), so \(A = \tfrac{1}{2}\). Setting \(x = \tfrac{1}{2}\) gives \(\tfrac{1}{4} = \tfrac{5}{4}B\), so \(B = \tfrac{1}{5}\). Setting \(x = -2\) gives \(-1 = 10C\), so \(C = -\tfrac{1}{10}\). Each fragment is a logarithm; for the middle one, the inner substitution \(u = 2x - 1\) contributes the factor \(\tfrac{1}{2}\):

\[ \int \frac{x^{2} + 2x - 1}{2x^{3} + 3x^{2} - 2x}\,dx = \frac{1}{2}\ln\lvert x\rvert + \frac{1}{10}\ln\lvert 2x - 1\rvert - \frac{1}{10}\ln\lvert x + 2\rvert + K, \]

on any interval avoiding \(0\), \(\tfrac{1}{2}\), and \(-2\). \qedmark
\end{envsolution}
\end{workedexample}
One habit belongs here. Once the constants are found, pick a value of \(x\) that has not been used yet, substitute it into the identity obtained after clearing denominators, and check that the two sides agree. In Example \ref{ex:8.18}, taking \(x = 1\) gives \(2\) on the left and \(3A + 3B + C = \tfrac{3}{2} + \tfrac{3}{5} - \tfrac{1}{10} = 2\) on the right. The check costs one line and catches sign and arithmetic slips before any integration begins, which matters because a wrong decomposition integrates just as smoothly as a correct one.

\begin{workedexample}
\begin{envexample}{Example}{ex:8.19}{}
Find \(\displaystyle\int \frac{x^{4} - 2x^{2} + 4x + 1}{x^{3} - x^{2} - x + 1}\,dx\).
\end{envexample}
\begin{envsolution}{Solution}{}{}
Improper again, so divide: the quotient is \(x + 1\) with remainder \(4x\), so

\[ \frac{x^{4} - 2x^{2} + 4x + 1}{x^{3} - x^{2} - x + 1} = x + 1 + \frac{4x}{x^{3} - x^{2} - x + 1}. \]

Factoring by grouping, \(x^{3} - x^{2} - x + 1 = x^{2}(x - 1) - (x - 1) = (x - 1)(x^{2} - 1) = (x - 1)^{2}(x + 1)\): a repeated linear factor. The template carries one fragment per power:

\[ \frac{4x}{(x - 1)^{2}(x + 1)} = \frac{A}{x - 1} + \frac{B}{(x - 1)^{2}} + \frac{C}{x + 1}. \]

Clearing denominators, \(4x = A(x - 1)(x + 1) + B(x + 1) + C(x - 1)^{2}\). Setting \(x = 1\) gives \(4 = 2B\), so \(B = 2\). Setting \(x = -1\) gives \(-4 = 4C\), so \(C = -1\). The cover-up shortcut cannot reach \(A\), as Strategy A.2 warns for repeated factors, so match the \(x^{2}\) coefficients: \(0 = A + C\), so \(A = 1\). Integrating fragment by fragment, with the power rule on the middle one,

\[ \begin{aligned}
          \int \frac{x^{4} - 2x^{2} + 4x + 1}{x^{3} - x^{2} - x + 1}\,dx &= \frac{x^{2}}{2} + x + \ln\lvert x - 1\rvert \\
            &\quad - \frac{2}{x - 1} - \ln\lvert x + 1\rvert + K,
        \end{aligned} \]

on any interval avoiding \(\pm 1\). \qedmark
\end{envsolution}
\end{workedexample}
\begin{envcaution}{Caution}{}{}
The template needs a fragment for every power of a repeated factor, not only the highest. Dropping the lower ones does not save work; it leaves a system with no solution at all. Writing \(\frac{4x}{(x - 1)^{2}(x + 1)}\) as \(\frac{B}{(x - 1)^{2}} + \frac{C}{x + 1}\) and clearing denominators gives \(4x = B(x + 1) + C(x - 1)^{2}\). Matching \(x^{2}\) coefficients forces \(C = 0\), matching \(x\) coefficients then forces \(B = 4\), and matching constants requires \(B + C = 0\). No such constants exist. The failure looks like an arithmetic mistake, but it is a sign that the template itself was wrong.
\end{envcaution}
\begin{workedexample}
\begin{envexample}{Example}{ex:8.20}{}
Find \(\displaystyle\int \frac{2x^{2} - x + 4}{x^{3} + 4x}\,dx\).
\end{envexample}
\begin{envsolution}{Solution}{}{}
The denominator factors as \(x(x^{2} + 4)\). The quadratic has \(p = 0\) and \(q = 4\), so \(p^{2} < 4q\) and it is irreducible; by Proposition A.7 it carries a linear numerator:

\[ \frac{2x^{2} - x + 4}{x(x^{2} + 4)} = \frac{A}{x} + \frac{Bx + C}{x^{2} + 4}. \]

Clearing denominators, \(2x^{2} - x + 4 = A(x^{2} + 4) + (Bx + C)x\). The constant term gives \(4 = 4A\), so \(A = 1\); the \(x^{2}\) coefficient gives \(2 = A + B\), so \(B = 1\); the \(x\) coefficient gives \(C = -1\). The quadratic fragment splits into its two standard pieces:

\[ \begin{aligned}
          \int \frac{x - 1}{x^{2} + 4}\,dx &= \frac{1}{2}\int \frac{2x}{x^{2} + 4}\,dx - \int \frac{dx}{x^{2} + 4} \\[2pt]
            &= \frac{1}{2}\ln\bigl(x^{2} + 4\bigr) - \frac{1}{2}\arctan\frac{x}{2} + K,
        \end{aligned} \]

the first by the substitution \(u = x^{2} + 4\) and the second by (8.A) with \(a = 2\). No absolute value is needed on \(x^{2} + 4\), which is positive for every \(x\). Assembling,

\[ \int \frac{2x^{2} - x + 4}{x^{3} + 4x}\,dx = \ln\lvert x\rvert + \frac{1}{2}\ln\bigl(x^{2} + 4\bigr) - \frac{1}{2}\arctan\frac{x}{2} + K, \]

on any interval avoiding \(0\). \qedmark
\end{envsolution}
\end{workedexample}
\begin{envcaution}{Caution}{}{}
Check irreducibility before assigning the \(Bx + C\) template: \(x^{2} + px + q\) is irreducible exactly when \(p^{2} < 4q\). A quadratic that looks unfactorable often factors after all. The factor \(x^{2} - x - 6\) splits as \((x - 3)(x + 2)\), and treating it as irreducible puts the whole decomposition on the wrong template. When in doubt, compute the discriminant first—\(p^{2} - 4q\) for the monic form above; for a general \(ax^{2} + bx + c\), factor out \(a\) first or use \(b^{2} - 4ac\), as Example \ref{ex:8.21} does.
\end{envcaution}
A repeated irreducible quadratic, the last case of Proposition A.7, adds no new integration idea: each power \(\frac{Bx + C}{(x^{2} + px + q)^{k}}\) again splits into a derivative-of-the-quadratic piece and a constant piece, and the constant pieces are handled by completing the square and a short reduction. The bookkeeping is longer, the shapes are the same, and this book leaves that case as stated.

One tool from step 4 has not yet appeared in action: completing the square.

\begin{workedexample}
\begin{envexample}{Example}{ex:8.21}{}
Find \(\displaystyle\int \frac{4x^{2} - 3x + 2}{4x^{2} - 4x + 3}\,dx\).
\end{envexample}
\begin{envsolution}{Solution}{}{}
The degrees are equal, so this fraction is improper and step 1 applies. Dividing,

\[ \frac{4x^{2} - 3x + 2}{4x^{2} - 4x + 3} = 1 + \frac{x - 1}{4x^{2} - 4x + 3}, \]

and the denominator has discriminant \((-4)^{2} - 4\cdot 4\cdot 3 < 0\), so it is irreducible and is its own decomposition. Completing the square, as in §A.5’s closing paragraph, \(4x^{2} - 4x + 3 = (2x - 1)^{2} + 2\). Substitute \(u = 2x - 1\), so \(x = \tfrac{u + 1}{2}\), \(dx = \tfrac{du}{2}\), and \(x - 1 = \tfrac{u - 1}{2}\):

\[ \begin{aligned}
          \int \frac{x - 1}{4x^{2} - 4x + 3}\,dx &= \frac{1}{4}\int \frac{u - 1}{u^{2} + 2}\,du \\[2pt]
            &= \frac{1}{8}\ln\bigl(u^{2} + 2\bigr) - \frac{1}{4}\cdot\frac{1}{\sqrt{2}}\arctan\frac{u}{\sqrt{2}} + C,
        \end{aligned} \]

splitting as in Example \ref{ex:8.20} and applying (8.A) with \(a = \sqrt{2}\). Returning to \(x\), with \(u^{2} + 2 = 4x^{2} - 4x + 3\) positive everywhere,

\[ \begin{aligned}
          \int \frac{4x^{2} - 3x + 2}{4x^{2} - 4x + 3}\,dx &= x + \frac{1}{8}\ln\bigl(4x^{2} - 4x + 3\bigr) \\
            &\quad - \frac{1}{4\sqrt{2}}\arctan\frac{2x - 1}{\sqrt{2}} + C,
        \end{aligned} \]

valid on all of \(\mathbb{R}\). This is the integral §A.5 promised: completing the square is the step that brings a general quadratic denominator within reach of (8.A). \qedmark
\end{envsolution}
\end{workedexample}
\subsechead{Rationalizing substitutions}
Some integrands are not rational but can be made so. When a single radical \(\sqrt[n]{g(x)}\) is the obstacle, the substitution \(u = \sqrt[n]{g(x)}\) often converts the whole integrand into a rational function of \(u\), where this section’s pipeline takes over.

\begin{workedexample}
\begin{envexample}{Example}{ex:8.22}{}
Find \(\displaystyle\int \frac{\sqrt{x + 4}}{x}\,dx\).
\end{envexample}
\begin{envsolution}{Solution}{}{}
Substitute \(u = \sqrt{x + 4}\), working on the component \(x > 0\), where \(u > 2\). Then \(u^{2} = x + 4\), so \(x = u^{2} - 4\) and \(dx = 2u\,du\):

\[ \int \frac{\sqrt{x + 4}}{x}\,dx = \int \frac{u}{u^{2} - 4}\,2u\,du = 2\int \frac{u^{2}}{u^{2} - 4}\,du. \]

The integrand is rational and improper, so divide: \(\frac{u^{2}}{u^{2} - 4} = 1 + \frac{4}{u^{2} - 4}\), and \(u^{2} - 4 = (u - 2)(u + 2)\) decomposes by cover-up as \(\frac{4}{(u-2)(u+2)} = \frac{1}{u - 2} - \frac{1}{u + 2}\). So

\[ 2\int \frac{u^{2}}{u^{2} - 4}\,du = 2u + 2\ln\lvert u - 2\rvert - 2\ln\lvert u + 2\rvert + C, \]

and returning to \(x\),

\[ \int \frac{\sqrt{x + 4}}{x}\,dx = 2\sqrt{x + 4} + 2\ln\left\lvert \frac{\sqrt{x + 4} - 2}{\sqrt{x + 4} + 2} \right\rvert + C. \]

Differentiating confirms the formula, and the same check shows it also serves on the other component, \(-4 < x < 0\). The solution chained three techniques in order: a substitution to remove the radical, a division to make the fraction proper, and partial fractions to finish. That chaining is how integration usually proceeds in practice, and the next section is about choosing such chains deliberately. \qedmark
\end{envsolution}
\end{workedexample}
Rational functions are the best-behaved integrands in this chapter: Strategy \ref{strat:8.5}’s four steps (divide, factor, decompose, integrate) always terminate in every case this section carries out, and the only borrowed ingredients are the factorization fact taken on credit above and the single fragment integral left as stated for repeated irreducible quadratic factors. Linear factors end in logarithms, their powers in power-rule terms, and irreducible quadratics in a logarithm plus an arctangent through (8.A). With parts, trigonometric methods, substitution, and partial fractions all in hand, the toolkit is complete. What remains is judgment: given an unfamiliar integral, which tool comes first? That question has a strategy of its own.


\sechead{8.5}{Strategy for Integration}
Each section so far answered the question “how does this technique work”. Practice asks a different question. An integral arrives unlabeled, and the work is to decide which tool fits, before any tool is applied. The four integrals

\[ \int \frac{\tan^{3}x}{\cos^{3}x}\,dx, \qquad \int e^{\sqrt{x}}\,dx, \qquad \int \frac{dx}{x\sqrt{\ln x}}, \qquad \int \sqrt{\frac{1 - x}{1 + x}}\,dx \]

look like near relatives: each is a small expression with one awkward feature. They go to four different tools, and none of them is best solved by the tool its form suggests first. Reading form is a skill of its own. It takes only seconds, a wrong guess costs little, and this section trains it.

\begin{envstrategy}{Strategy}{strat:8.6}{A four-step strategy}
\begin{steps}
  \item \emph{Simplify the integrand first.} Algebra and identities can change the problem’s class entirely: expand or split fractions, rationalize, apply a trigonometric identity. Only then judge what kind of integral it is.
  \item \emph{Look for an obvious substitution.} If some inner function \(g(x)\) appears together with its derivative, up to a constant factor, take \(u = g(x)\) (Strategy 6.2) and finish before reaching for anything heavier.
  \item \emph{Classify by form and dispatch.} Powers or products of trigonometric functions go to Strategies \ref{strat:8.2} and \ref{strat:8.3} (both in §8.2). A radical over a quadratic goes to Strategy \ref{strat:8.4} (§8.3), completing the square first when the quadratic is not centered. A rational function goes to Strategy \ref{strat:8.5}. A product of unrelated factors, or a lone logarithm or inverse function, goes to parts by Strategy \ref{strat:8.1}.
  \item \emph{If nothing has worked, try again from another angle.} Manipulate the integrand into a relative of a known integral, substitute to expose hidden structure, or chain techniques in sequence, as Examples \ref{ex:8.22} and \ref{ex:8.24} do. Several routes may succeed on the same integral; a dead end costs only the time it took.
\end{steps}
\end{envstrategy}
Step 3 leans on a stock of antiderivatives that can be recognized on sight. One scaled form joins the list first: differentiating \(\arcsin\tfrac{x}{a}\) by the chain rule (Theorem 3.3) and Example 3.14 gives \(\tfrac{1}{\sqrt{a^{2} - x^{2}}}\), so \(\int\tfrac{dx}{\sqrt{a^{2} - x^{2}}} = \arcsin\tfrac{x}{a} + C\) for \(a > 0\) on \((-a, a)\). The working table, with each entry’s source in parentheses:

\[ \begin{aligned}
    &\int x^{n}\,dx = \frac{x^{n+1}}{n+1} + C \;(n \ne -1), \qquad \int \frac{dx}{x} = \ln\lvert x\rvert + C \quad(\text{§6.4}), \\[2pt]
    &\int e^{x}\,dx = e^{x} + C, \qquad \int \sin x\,dx = -\cos x + C \quad(\text{§6.4}), \\[2pt]
    &\int \cos x\,dx = \sin x + C, \qquad \int \sec^{2}x\,dx = \tan x + C \quad(\text{§6.4}), \\[2pt]
    &\int \tan x\,dx = \ln\lvert\sec x\rvert + C \quad(\text{Prop.\ 8.1}), \\[2pt]
    &\int \sec x\,dx = \ln\lvert\sec x + \tan x\rvert + C \quad(\text{Prop.\ 8.1}), \\[2pt]
    &\int \frac{dx}{x^{2} + a^{2}} = \frac{1}{a}\arctan\frac{x}{a} + C \quad(\text{8.A}), \\[2pt]
    &\int \frac{dx}{\sqrt{a^{2} - x^{2}}} = \arcsin\frac{x}{a} + C \quad(\text{above}).
  \end{aligned} \]

We now run the opening’s four integrals through the strategy, in order.

\begin{workedexample}
\begin{envexample}{Example}{ex:8.23}{}
Find \(\displaystyle\int \frac{\tan^{3}x}{\cos^{3}x}\,dx\).
\end{envexample}
\begin{envsolution}{Solution}{}{}
Step 1: since \(\tfrac{1}{\cos x} = \sec x\), the integrand is \(\tan^{3}x\,\sec^{3}x\), a tangent–secant power. Step 3 dispatches it to Strategy \ref{strat:8.3}: the power of tangent is odd, so save a factor of \(\sec x\tan x\) and convert the remaining tangents by \(\tan^{2}x = \sec^{2}x - 1\):

\[ \int \tan^{3}x\,\sec^{3}x\,dx = \int \bigl(\sec^{2}x - 1\bigr)\sec^{2}x\;\sec x\tan x\,dx. \]

With \(u = \sec x\), \(du = \sec x\tan x\,dx\), this is \(\int(u^{4} - u^{2})\,du\), so

\[ \int \frac{\tan^{3}x}{\cos^{3}x}\,dx = \frac{\sec^{5}x}{5} - \frac{\sec^{3}x}{3} + C. \]

The only decision was the rewrite in step 1; after it, the machinery of Strategy \ref{strat:8.3} finished the integral with no further choices. \qedmark
\end{envsolution}
\end{workedexample}
\begin{workedexample}
\begin{envexample}{Example}{ex:8.24}{}
Find \(\displaystyle\int e^{\sqrt{x}}\,dx\) for \(x > 0\).
\end{envexample}
\begin{envsolution}{Solution}{}{}
No product, no trigonometry, no rational structure. The awkward feature is the radical in the exponent, and step 4’s advice is to substitute it away even though \(du\) is nowhere in sight. Let \(u = \sqrt{x}\), so \(x = u^{2}\) and \(dx = 2u\,du\):

\[ \int e^{\sqrt{x}}\,dx = \int e^{u}\,2u\,du = 2\int u\,e^{u}\,du. \]

The substitution has manufactured a product of unrelated factors, which is parts territory. As in Example \ref{ex:8.3}, \(\int u e^{u}\,du = (u - 1)e^{u} + C\), so

\[ \int e^{\sqrt{x}}\,dx = 2\bigl(\sqrt{x} - 1\bigr)e^{\sqrt{x}} + C. \]

Two techniques chained: a substitution to normalize the exponent, then parts to clear the leftover factor. Neither alone finishes the integral. \qedmark
\end{envsolution}
\end{workedexample}
\begin{workedexample}
\begin{envexample}{Example}{ex:8.25}{}
Find \(\displaystyle\int \frac{dx}{x\sqrt{\ln x}}\) for \(x > 1\).
\end{envexample}
\begin{envsolution}{Solution}{}{}
Step 2 settles this one before any classification: the factor \(\tfrac{dx}{x}\) is exactly the derivative of \(\ln x\), so \(u = \ln x\) is the obvious substitution, with \(u > 0\) on the stated domain:

\[ \int \frac{dx}{x\sqrt{\ln x}} = \int u^{-1/2}\,du = 2\sqrt{u} + C = 2\sqrt{\ln x} + C. \]

An integrand that looks exotic can have a two-line solution. Checking for the visible \(du\) first is what keeps the heavier tools in reserve. \qedmark
\end{envsolution}
\end{workedexample}
\begin{workedexample}
\begin{envexample}{Example}{ex:8.26}{}
Find \(\displaystyle\int \sqrt{\frac{1 - x}{1 + x}}\,dx\) on \((-1, 1)\).
\end{envexample}
\begin{envsolution}{Solution}{}{}
A rationalizing substitution in the style of Example \ref{ex:8.22} would work, but step 1 offers something faster. Multiply inside the root by \(\tfrac{1 - x}{1 - x}\):

\[ \sqrt{\frac{1 - x}{1 + x}} = \sqrt{\frac{(1 - x)^{2}}{1 - x^{2}}} = \frac{1 - x}{\sqrt{1 - x^{2}}}, \]

where dropping the absolute value is justified because \(1 - x > 0\) on \((-1, 1)\). The integrand splits into two table entries:

\[ \int \frac{dx}{\sqrt{1 - x^{2}}} - \int \frac{x}{\sqrt{1 - x^{2}}}\,dx = \arcsin x + \sqrt{1 - x^{2}} + C, \]

the first from the table above with \(a = 1\), the second by the substitution \(w = 1 - x^{2}\), \(dw = -2x\,dx\). One algebraic move changed the problem’s class from “radical of a ratio” to “two standard forms”. That is what step 1 is for. \qedmark
\end{envsolution}
\end{workedexample}
Two remarks are worth adding here, one about step 1 and one about step 4. First, an algebraic rewriting need not hold on the whole domain. Dropping an absolute value, cancelling a factor, or extracting a square root can each change the set of \(x\) on which the two sides agree, and can do so silently, so name the interval at the moment you rewrite, as was done in Example \ref{ex:8.26}.

Second, when several routes succeed, their answers need not look alike. Two antiderivatives of the same function on the same interval differ by a constant (Corollary 4.4), and by nothing more, so an answer that looks different from another may still be correct. There are two ways to settle the question: differentiate and compare with the integrand, or subtract the two answers and simplify. This chapter already contains two such cases. In Example \ref{ex:8.8} the power-reduction answer looks unlike the one the reduction formula of Example \ref{ex:8.6} gives, and the two agree only after a double-angle identity; in Example \ref{ex:8.15} the two forms differ by the constant \(\ln a\).

\subsechead{The limits of the toolkit}
A fair question after five sections of techniques: does every integrand now yield to the toolkit? Try the innocent-looking \(\int e^{-x^{2}}\,dx\). It is not a product, no substitution exposes an inner derivative, it is neither trigonometric nor rational, and no rewriting in this chapter changes that. None of the techniques of this book produces an antiderivative for it that is built from the functions we have named. Whether \emph{any} formula of that kind exists is a question this book does not take up.

What matters here is practical, and it is twofold. First, antiderivatives are not the only way to a definite integral: the next section shows that \(\int_{1}^{\infty} e^{-x^{2}}\,dx\) is a perfectly well-defined finite number even without a formula for its antiderivative. Second, when a number is what we need, §8.7’s rules compute it to any stated accuracy. The toolkit’s boundary is real, and the chapter’s last two sections are what lie on the far side of it.

There are further systematic conversions beyond this book’s needs, including one that turns any rational expression in \(\sin x\) and \(\cos x\) into an ordinary rational function; the four steps above are enough for everything this book asks.

The strategy is short because the work was done in the sections it points to: simplify, check for a visible substitution, classify, and dispatch to the section that covers the form, chaining techniques when the first leaves a new form for the second. The four examples were each decided by the step named in their solutions, not by new machinery. What follows next is not another technique but a widening of the integral itself, to infinite intervals and unbounded integrands, where the first question is no longer “which tool” but “does this integral mean anything at all”.


\sechead{8.6}{Improper Integrals}
Every integral so far has lived on a closed interval \([a, b]\): Definition 6.2 builds its Riemann sums from finitely many slices of a bounded interval, and for the continuous integrands this book evaluates, Theorem 6.1 guarantees the sums converge. Two natural questions leave that setting. What is the area under \(y = 1/x^{2}\) from \(x = 1\) out to infinity? And what is the area under \(y = 1/\sqrt{x - 2}\) from \(2\) to \(5\), where the integrand blows up at the left end? In both cases the symbol \(\int\) has, as it stands, no meaning: no Riemann sum machinery applies. But the book already owns a device for questions like these. Truncate to a proper integral, and ask what the truncations approach. For the first region, \(\int_{1}^{t} x^{-2}\,dx = 1 - \tfrac{1}{t}\), and as \(t\) grows these values settle toward \(1\). An endless region can hold a finite area, and a limit is how we say so. This section extends the integral by exactly that move, made twice: once for infinite intervals, once for unbounded integrands.

\subsechead{Type 1: infinite intervals}
\begin{envdefinition}{Definition}{def:8.1}{Improper integral of Type 1}
(a) If \(f\) is continuous on \([a, \infty)\), then

\[ \int_{a}^{\infty} f(x)\,dx = \lim_{t \to \infty} \int_{a}^{t} f(x)\,dx, \]

provided the limit exists as a finite number; the improper integral is then called \emph{convergent}, and \emph{divergent} otherwise.

(b) If \(f\) is continuous on \((-\infty, b]\), then \(\displaystyle\int_{-\infty}^{b} f(x)\,dx = \lim_{t \to -\infty} \int_{t}^{b} f(x)\,dx\), with the same vocabulary.

(c) If \(f\) is continuous on \((-\infty, \infty)\), choose any number \(c\) and set

\[ \int_{-\infty}^{\infty} f(x)\,dx = \int_{-\infty}^{c} f(x)\,dx + \int_{c}^{\infty} f(x)\,dx, \]

provided \emph{both} integrals on the right converge.
\end{envdefinition}
The choice of \(c\) in part (c) does not matter: additivity (Theorem 6.2, part 3) shifts any finite stretch between the two halves without changing their sum, so one split point gives a convergent sum exactly when another does, with the same value. Note also what part (c) does \emph{not} say; see the caution below.

\begin{workedexample}
\begin{envexample}{Example}{ex:8.27}{}
Decide whether each integral converges, and evaluate the convergent one:

\[ \text{(a)}\ \int_{1}^{\infty} \frac{dx}{x}, \qquad \text{(b)}\ \int_{1}^{\infty} \frac{dx}{x^{2}}. \]
\end{envexample}
\begin{envsolution}{Solution}{}{}
(a) For each \(t > 1\),

\[ \int_{1}^{t} \frac{dx}{x} = \Bigl[\ln x\Bigr]_{1}^{t} = \ln t, \]

and \(\ln t\) grows without bound as \(t \to \infty\). The limit is not a finite number, so \(\int_{1}^{\infty} \tfrac{dx}{x}\) \emph{diverges}.

(b) For each \(t > 1\),

\[ \int_{1}^{t} \frac{dx}{x^{2}} = \Bigl[-\frac{1}{x}\Bigr]_{1}^{t} = 1 - \frac{1}{t} \longrightarrow 1, \]

so \(\int_{1}^{\infty} \tfrac{dx}{x^{2}}\) \emph{converges}, with value \(1\).

The two curves look alike: both positive, both decreasing to \(0\). The verdicts differ because the tails thin at different rates. The values of \(1/x^{2}\) decay fast enough that the accumulated area stalls at \(1\); the values of \(1/x\) do not, and the area eventually passes any bound. How fast the integrand decays is the entire question, and the \(p\)-test below makes the boundary exact. \qedmark
\end{envsolution}
\end{workedexample}
\begin{figureblock}
  \includegraphics[width=122.22mm]{ch08/figs/tail-comparison}
\figcaption{fig:8.6}{The tails of \(y = 1/x\) and \(y = 1/x^{2}\) on \([1, \infty)\). The curves meet only at \((1, 1)\) and both decay to \(0\); the darker region under \(1/x^{2}\) stalls at total area \(1\), while the region under \(1/x\) keeps growing as the cutoff \(t\) moves right. How fast the tail thins decides the verdict.}
\end{figureblock}
\begin{workedexample}
\begin{envexample}{Example}{ex:8.28}{}
Evaluate \(\displaystyle\int_{-\infty}^{\infty} \frac{dx}{1 + x^{2}}\).
\end{envexample}
\begin{envsolution}{Solution}{}{}
By Definition \ref{def:8.1}(c), split at \(c = 0\) and treat each half. For the right half,

\[ \int_{0}^{t} \frac{dx}{1 + x^{2}} = \Bigl[\arctan x\Bigr]_{0}^{t} = \arctan t \longrightarrow \frac{\pi}{2}, \]

since \(y = \tfrac{\pi}{2}\) is the horizontal asymptote of the arctangent (§1.2). By the symmetric computation, \(\int_{-t}^{0} \to \tfrac{\pi}{2}\) as well. Both halves converge, so the integral converges and

\[ \int_{-\infty}^{\infty} \frac{dx}{1 + x^{2}} = \frac{\pi}{2} + \frac{\pi}{2} = \pi. \]

An unbounded region whose area is exactly \(\pi\): the curve decays like \(1/x^{2}\) in both directions, fast enough for the \(p\)-test’s convergent side. \qedmark
\end{envsolution}
\end{workedexample}
\begin{envcaution}{Caution}{}{}
Definition \ref{def:8.1}(c) requires both halves to converge \emph{separately}. It is not the same as \(\lim_{t \to \infty}\int_{-t}^{t} f(x)\,dx\), and the difference matters. For \(f(x) = x\), every symmetric truncation gives \(\int_{-t}^{t} x\,dx = 0\), yet each half of \(\int_{-\infty}^{\infty} x\,dx\) diverges, so the integral diverges. Symmetric cancellation of two infinite areas is not convergence.
\end{envcaution}
\begin{envcaution}{Caution}{}{}
Example \ref{ex:8.27}(a) diverged because its truncations grew without bound, and that is the most familiar way for the limit to fail. It is not the only way. \emph{Divergent} says that the limit is not a finite number, which also covers truncations that never settle. For \(f(x) = \cos x\), the truncations are \(\int_{0}^{t}\cos x\,dx = \sin t\), which returns to \(1\) and to \(-1\) again and again as \(t\) grows, so no single number is being approached and \(\int_{0}^{\infty}\cos x\,dx\) diverges. Divergence is the absence of a finite limit, not a claim that the area is infinite.
\end{envcaution}
\begin{workedexample}
\begin{envexample}{Example}{ex:8.29}{}
Evaluate \(\displaystyle\int_{0}^{\infty} t e^{-t}\,dt\).
\end{envexample}
\begin{envsolution}{Solution}{}{}
On each interval \([0, s]\), integrate by parts (Theorem \ref{thm:8.2}) with \(u = t\), \(dv = e^{-t}\,dt\), so \(du = dt\) and \(v = -e^{-t}\):

\[ \int_{0}^{s} t e^{-t}\,dt = \Bigl[-t e^{-t}\Bigr]_{0}^{s} + \int_{0}^{s} e^{-t}\,dt = -s e^{-s} + \bigl(1 - e^{-s}\bigr). \]

As \(s \to \infty\), the exponential term \(e^{-s} \to 0\), and \(s e^{-s} = s/e^{s}\) is an \(\infty/\infty\) form that L’Hôpital’s Rule (Theorem 5.5) settles: the quotient has the same limit as \(1/e^{s}\), namely \(0\). So the truncations approach \(0 + 1 - 0\), and

\[ \int_{0}^{\infty} t e^{-t}\,dt = 1. \]

The techniques and the extended definition worked one after the other: parts produced the truncated value, and the limit made sense of the infinite interval. \qedmark
\end{envsolution}
\end{workedexample}
The convergence recorded in Example \ref{ex:8.27}(b) has physical content. Newton’s inverse-square law, an empirical model in the same sense as Hooke’s law in §7.4, says that a body of mass \(M\) pulls a mass \(m\) at distance \(x\) with force \(GMm/x^{2}\). By Definition 7.3, the work needed to move the mass from the body’s surface, at distance \(R\) from its centre, out to distance \(t\) is

\[ \int_{R}^{t} \frac{GMm}{x^{2}}\,dx = GMm\Bigl(\frac{1}{R} - \frac{1}{t}\Bigr) \longrightarrow \frac{GMm}{R}, \]

so the improper integral converges: moving the mass arbitrarily far from the body requires only a finite amount of work. The reason is the one Example \ref{ex:8.27}(b) recorded: the exponent \(2\) sits on the convergent side of the boundary.

Example \ref{ex:8.27}’s two verdicts can be produced by one curve: rotating \(y = 1/x\) on \([1, t]\) about the \(x\)-axis gives a solid, traditionally called Gabriel’s horn, whose volume \(\int_{1}^{t}\pi x^{-2}\,dx\) approaches \(\pi\), while its lateral surface area is at least \(\int_{1}^{t} 2\pi x^{-1}\,dx = 2\pi\ln t\) and grows without bound as \(t\) does, since the Theorem 7.4 integrand \(2\pi x^{-1}\sqrt{1 + x^{-4}}\) is never smaller than \(2\pi x^{-1}\) (Theorem 6.2, part 4).

\subsechead{Type 2: unbounded integrands}
The second violation keeps the interval finite but lets the integrand blow up at an endpoint. The extension move is the same: truncate short of the bad point, and take a limit.

\begin{envdefinition}{Definition}{def:8.2}{Improper integral of Type 2}
(a) If \(f\) is continuous on \([a, b)\) and unbounded near \(b\), then

\[ \int_{a}^{b} f(x)\,dx = \lim_{t \to b^{-}} \int_{a}^{t} f(x)\,dx, \]

provided the limit exists as a finite number; the integral is then \emph{convergent}, and \emph{divergent} otherwise.

(b) If \(f\) is continuous on \((a, b]\) and unbounded near \(a\), then \(\displaystyle\int_{a}^{b} f(x)\,dx = \lim_{t \to a^{+}} \int_{t}^{b} f(x)\,dx\), with the same vocabulary.

(c) If \(f\) is unbounded near an interior point \(c\) of \([a, b]\) and continuous elsewhere on the interval, split there: \(\displaystyle\int_{a}^{b} f(x)\,dx = \int_{a}^{c} f(x)\,dx + \int_{c}^{b} f(x)\,dx\), provided \emph{both} halves converge.
\end{envdefinition}
\begin{workedexample}
\begin{envexample}{Example}{ex:8.30}{}
Decide whether each integral converges, and evaluate where possible:

\[ \text{(a)}\ \int_{2}^{5} \frac{dx}{\sqrt{x - 2}}, \qquad \text{(b)}\ \int_{0}^{1} \frac{dx}{x^{p}} \ \text{ for } p > 0. \]
\end{envexample}
\begin{envsolution}{Solution}{}{}
(a) The integrand is continuous on \((2, 5]\) and unbounded near \(2\), so Definition \ref{def:8.2}(b) applies. For \(2 < t < 5\),

\[ \int_{t}^{5} \frac{dx}{\sqrt{x - 2}} = \Bigl[2\sqrt{x - 2}\Bigr]_{t}^{5} = 2\sqrt{3} - 2\sqrt{t - 2} \longrightarrow 2\sqrt{3} \]

as \(t \to 2^{+}\). This answers the opener’s second question: the integral converges to \(2\sqrt{3}\), because the spike at \(x = 2\) is tall but thin enough.

(b) For \(p \ne 1\) and \(0 < t < 1\),

\[ \int_{t}^{1} x^{-p}\,dx = \frac{1 - t^{1 - p}}{1 - p}. \]

If \(p < 1\), then \(t^{1-p} \to 0\) as \(t \to 0^{+}\), and the integral converges to \(\tfrac{1}{1 - p}\). If \(p > 1\), then \(t^{1-p}\) grows without bound and the integral diverges; for \(p = 1\), \(\int_{t}^{1} \tfrac{dx}{x} = -\ln t\) also grows without bound. So \(\int_{0}^{1} x^{-p}\,dx\) converges exactly when \(p < 1\). This is the mirror image of the \(p\)-test proved below: at \(0\) the mild powers \((p < 1)\) converge, at \(\infty\) the steep ones \((p > 1)\) do, and \(p = 1\) fails at both ends. \qedmark
\end{envsolution}
\end{workedexample}
\begin{figureblock}
  \includegraphics[width=114.29mm]{ch08/figs/type2-truncation}
\figcaption{fig:8.7}{The truncation of Definition \ref{def:8.2}(b) for \(\int_{2}^{5} dx/\sqrt{x - 2}\): the region from \(t\) to \(5\) is a proper integral, and as \(t \to 2^{+}\) the limit captures the lighter sliver against the asymptote. The spike at \(x = 2\) is tall but thin enough for the limit to stay finite.}
\end{figureblock}
\begin{workedexample}
\begin{envexample}{Example}{ex:8.31}{}
What is wrong with the computation \(\displaystyle\int_{-1}^{3} \frac{dx}{x^{2}} = \Bigl[-\frac{1}{x}\Bigr]_{-1}^{3} = -\frac{1}{3} - 1 = -\frac{4}{3}\)?
\end{envexample}
\begin{envsolution}{Solution}{}{}
The answer refutes itself: the integrand is positive, and an integral that adds up positive values can never come out negative. The error is upstream. The Fundamental Theorem applies to integrands continuous on the whole interval, and \(1/x^{2}\) is unbounded near the interior point \(0\). The integral is improper, and Definition \ref{def:8.2}(c) requires splitting there:

\[ \int_{-1}^{3} \frac{dx}{x^{2}} = \int_{-1}^{0} \frac{dx}{x^{2}} + \int_{0}^{3} \frac{dx}{x^{2}}, \]

provided both halves converge. They do not: near \(0\) the integrand is a \(p = 2\) power, and by the computation of Example \ref{ex:8.30}(b), repeated on the interval \((0, 3]\), the right half diverges. So \(\int_{-1}^{3} x^{-2}\,dx\) diverges, and the number \(-\tfrac{4}{3}\) was meaningless. \qedmark
\end{envsolution}
\end{workedexample}
\begin{envcaution}{Caution}{}{}
Before applying the Fundamental Theorem on \([a, b]\), check the whole interval for points where the integrand is undefined or unbounded. A singularity in the interior does not announce itself in the limits of integration, and integrating across it, as in Example \ref{ex:8.31}, produces a definite-looking number with no meaning. If a singularity is found, split there and test each piece by Definition \ref{def:8.2}.
\end{envcaution}
\begin{figureblock}
  \includegraphics[width=119.58mm]{ch08/figs/interior-singularity}
\figcaption{fig:8.8}{The integrand of Example \ref{ex:8.31}: two branches of \(y = 1/x^{2}\) climbing the same vertical asymptote, which coincides with the \(y\)-axis. The limits \(-1\) and \(3\) sit harmlessly on the axis while the singularity waits between them. Scan the whole interval before applying the Fundamental Theorem.}
\end{figureblock}
\subsechead{The \(p\)-test and the Comparison Theorem}
Example \ref{ex:8.27}’s contrast has an exact boundary, and it is worth finding. The family \(\int_{1}^{\infty} x^{-p}\,dx\) converges or diverges according to the exponent, and the result is used often enough to record permanently.

\begin{envproposition}{Proposition}{prop:8.2}{The \(p\)-test}
The improper integral

\[ \int_{1}^{\infty} \frac{dx}{x^{p}} \]

converges if \(p > 1\), with value \(\dfrac{1}{p - 1}\), and diverges if \(p \le 1\).
\end{envproposition}
\begin{envproof}{Proof}{}{}
For \(p \ne 1\), the power rule gives, for \(t > 1\),

\[ \int_{1}^{t} x^{-p}\,dx = \frac{t^{1 - p} - 1}{1 - p}. \]

If \(p > 1\), the exponent \(1 - p\) is negative, so \(t^{1-p} \to 0\) and the truncations approach \(\tfrac{-1}{1 - p} = \tfrac{1}{p - 1}\): convergence, with the stated value. If \(p < 1\), the exponent is positive, \(t^{1-p}\) grows without bound, and the integral diverges. This already covers \(p \le 0\), where the integrand does not even tend to \(0\): there \(x^{-p} \ge 1\) on \([1, \infty)\), so \(\int_{1}^{t} x^{-p}\,dx \ge t - 1\) by comparison (Theorem 6.2, part 4), and divergence is direct. Finally, for \(p = 1\), Example \ref{ex:8.27}(a) showed the truncations equal \(\ln t\), which grows without bound. \qedmark
\end{envproof}
The last tool answers a question the \(p\)-test cannot. Section 8.5 met \(\int e^{-x^{2}}\,dx\), an integrand for which none of our techniques produces an antiderivative, so no truncation can be evaluated in closed form and the definition cannot be applied directly. Yet convergence is a question about \emph{size}, not about formulas. If a positive integrand sits below one whose improper integral converges, its accumulated area is trapped below a finite ceiling, and it should have nowhere to go but a limit. The next theorem makes that argument honest, and its proof is the section’s one genuinely theoretical moment: everything comes from machinery the book already owns.

\begin{envtheorem}{Theorem}{thm:8.3}{Comparison Theorem for Improper Integrals}
Let \(f\) and \(g\) be continuous with \(0 \le f(x) \le g(x)\) for \(x \ge a\).

(a) If \(\displaystyle\int_{a}^{\infty} g(x)\,dx\) converges, then \(\displaystyle\int_{a}^{\infty} f(x)\,dx\) converges, and its value is at most \(\displaystyle\int_{a}^{\infty} g(x)\,dx\).

(b) If \(\displaystyle\int_{a}^{\infty} f(x)\,dx\) diverges, then \(\displaystyle\int_{a}^{\infty} g(x)\,dx\) diverges.
\end{envtheorem}
\begin{envproof}{Proof}{}{}
(a) Write \(F(t) = \int_{a}^{t} f(x)\,dx\) and \(G(t) = \int_{a}^{t} g(x)\,dx\), and let \(M = \int_{a}^{\infty} g(x)\,dx\), the finite limit assumed to exist. Three observations set up the argument.

First, \(F\) is increasing: for \(t' > t\), additivity (Theorem 6.2, part 3) gives \(F(t') - F(t) = \int_{t}^{t'} f(x)\,dx \ge 0\), since \(f \ge 0\) (Theorem 6.2, part 4). The same reasoning makes \(G\) increasing.

Second, \(G(t) \le M\) for every \(t\). An increasing function cannot overshoot its own limit: if some \(G(t_{0})\) exceeded \(M\), then every later \(G(t)\) would be at least \(G(t_{0})\), and the values could not approach \(M\).

Third, \(F(t) \le G(t)\) for every \(t\), by comparison on \([a, t]\) (Theorem 6.2, part 4). Combining: \(F\) is increasing and \(F(t) \le M\) for all \(t \ge a\).

Now sample \(F\) along the integers, since Theorem 4.1’s monotone principle speaks of sequences, not of functions of a real variable. The numbers \(F(a + n)\), for \(n = 1, 2, 3, \dots\), form an increasing sequence bounded above by \(M\), so by Theorem 4.1 they converge to some limit \(L \le M\). Since an increasing sequence cannot overshoot its own limit, each \(F(a + n) \le L\).

Finally, the sampled limit controls the whole function. Given any \(t \ge a + 1\), let \(n\) be the greatest integer not exceeding \(t - a\), so that \(a + n \le t < a + n + 1\). Since \(F\) is increasing,

\[ F(a + n) \;\le\; F(t) \;\le\; F(a + n + 1) \;\le\; L. \]

As \(t \to \infty\), the index \(n\) grows with it, and \(F(a + n)\) approaches \(L\). Every value \(F(t)\) is thus squeezed between numbers that approach \(L\) and the ceiling \(L\) itself, so \(F(t) \to L\). By Definition \ref{def:8.1}, \(\int_{a}^{\infty} f(x)\,dx\) converges, with value \(L \le M\).

(b) This is the contrapositive of (a): if \(\int_{a}^{\infty} g\) converged, then (a) would force \(\int_{a}^{\infty} f\) to converge as well. \qedmark
\end{envproof}
\begin{figureblock}
  \includegraphics[width=122.22mm]{ch08/figs/monotone-sampling}
\figcaption{fig:8.9}{The proof of Theorem \ref{thm:8.3}(a) in one schematic picture: the increasing accumulation \(F\) stays below its ceiling \(M\), the integer samples \(F(a + n)\) climb to their limit \(L\) (Theorem 4.1), and every \(t\) is bracketed between consecutive samples, so \(F(t)\) is squeezed to \(L\) as well. The curve’s shape is schematic.}
\end{figureblock}
\begin{envcaution}{Caution}{}{}
Theorem \ref{thm:8.3} assumes \(0 \le f(x) \le g(x)\), and that assumption cannot be dropped. For an integrand that changes sign, a comparison of this kind yields nothing at all: the proof of Theorem \ref{thm:8.3} rests on \(F\) being increasing, which fails once \(f\) takes negative values, so neither convergence nor divergence follows. Such integrals are outside what this book takes up; every comparison in this chapter therefore begins by checking that both integrands are non-negative.
\end{envcaution}
\begin{workedexample}
\begin{envexample}{Example}{ex:8.32}{}
Show that \(\displaystyle\int_{1}^{\infty} e^{-x^{2}}\,dx\) converges.
\end{envexample}
\begin{envsolution}{Solution}{}{}
For \(x \ge 1\) we have \(x^{2} \ge x\), so \(-x^{2} \le -x\), and since the exponential function is increasing, \(0 < e^{-x^{2}} \le e^{-x}\). The larger integrand’s integral converges by direct computation:

\[ \int_{1}^{t} e^{-x}\,dx = e^{-1} - e^{-t} \longrightarrow \frac{1}{e}. \]

By the Comparison Theorem (Theorem \ref{thm:8.3}(a)), \(\int_{1}^{\infty} e^{-x^{2}}\,dx\) converges, with value at most \(\tfrac{1}{e}\). Over \([0, \infty)\) the conclusion extends at once, since \(\int_{0}^{1} e^{-x^{2}}\,dx\) is an ordinary proper integral: the full integral \(\int_{0}^{\infty} e^{-x^{2}}\,dx\) converges too.

This answers the question §8.5 left open. We have no formula for its antiderivative, yet the improper integral is a definite finite number, pinned down without ever being evaluated. Section 8.7 is about computing such numbers. \qedmark
\end{envsolution}
\end{workedexample}
\begin{figureblock}
  \includegraphics[width=117.35mm]{ch08/figs/comparison-trap}
\figcaption{fig:8.10}{The Comparison Theorem’s trap in Example \ref{ex:8.32}: on \([1, \infty)\) the graph of \(e^{-x^{2}}\) runs below \(e^{-x}\), whose tail holds the finite area \(1/e\), so the darker region is confined below a finite roof. The curves meet at \(x = 1\) and separate at a visible angle.}
\end{figureblock}
\begin{workedexample}
\begin{envexample}{Example}{ex:8.33}{}
Show that \(\displaystyle\int_{3}^{\infty} \frac{dx}{\sqrt{x^{2} - x}}\) diverges.
\end{envexample}
\begin{envsolution}{Solution}{}{}
First check the hypotheses. On \([3, \infty)\) we have \(x^{2} - x = x(x - 1) \ge 6\), so the integrand is defined, continuous, and positive there.

Divergence is the target, so by Theorem \ref{thm:8.3}(b) the integrand must be placed \emph{above} something already known to diverge. For \(x \ge 3\), \(0 < x^{2} - x < x^{2}\), and taking square roots gives \(\sqrt{x^{2} - x} < x\). Reciprocals reverse the inequality:

\[ 0 \;<\; \frac{1}{x} \;<\; \frac{1}{\sqrt{x^{2} - x}} \qquad (x \ge 3). \]

In the notation of Theorem \ref{thm:8.3} this makes \(f(x) = 1/x\) the lower function and \(g(x) = 1/\sqrt{x^{2} - x}\) the upper one, both continuous on \([3, \infty)\). The lower integral diverges: \(\int_{3}^{t}\tfrac{dx}{x} = \ln t - \ln 3\) by §6.4 and Theorem 6.4, and this grows without bound, so Definition \ref{def:8.1}(a) records divergence. This is Example \ref{ex:8.27}(a)’s computation from a different starting point.

Theorem \ref{thm:8.3}(b) therefore gives that \(\int_{3}^{\infty}\tfrac{dx}{\sqrt{x^{2} - x}}\) diverges.

Getting the direction right is the whole point here. It is just as easy to check that \(\tfrac{1}{\sqrt{x^{2}-x}} \le \tfrac{2}{x}\) on \([3, \infty)\): after squaring and clearing denominators, this is the inequality \(3x^{2} \ge 4x\). The integral \(\int_{3}^{\infty}\tfrac{2\,dx}{x}\) diverges as well. That comparison proves nothing: a divergent upper bound carries no information. To establish divergence, the known-divergent integral has to be the one underneath.

This integrand does have an antiderivative, reachable by completing the square and a secant substitution (§8.3). Deciding convergence did not require it. \qedmark
\end{envsolution}
\end{workedexample}
\begin{envcaution}{Caution}{}{}
Even when \(0 \le f \le g\) does hold, only two of the four ways of pairing a function with a verdict give a conclusion, and Example \ref{ex:8.33} has just used one of those two. The other two give nothing: if \(\int_{a}^{\infty} f\) converges, no conclusion follows about \(\int_{a}^{\infty} g\); if \(\int_{a}^{\infty} g\) diverges, none follows about \(\int_{a}^{\infty} f\). Example \ref{ex:8.27} supplies a counterexample to both at once: on \([1, \infty)\) we have \(1/x^{2} \le 1/x\), the smaller integral converges, and the larger one diverges. Decide first which verdict you are trying to prove, then look for a comparison function on the matching side.
\end{envcaution}
One consequence of additivity deserves stating on its own: convergence is decided by the tail. Attaching or removing a proper integral over a finite interval changes the value but not the verdict (Theorem 6.2, part 3). The practical use is that a comparison need not hold from the very start. If \(0 \le f \le g\) only for \(x \ge b\), with \(b > a\), split \(\int_{a}^{\infty} f\) at \(b\), apply Theorem \ref{thm:8.3} to \(\int_{b}^{\infty} f\), and add back \(\int_{a}^{b} f\), which is an ordinary proper integral whenever \(f\) is continuous on \([a, b]\) (Theorem 6.1). Example \ref{ex:8.32} did exactly this: its comparison \(e^{-x^{2}} \le e^{-x}\) holds only for \(x \ge 1\), so the verdict obtained on \([1, \infty)\) carried across to \([0, \infty)\).

Mixed cases reduce to the two types. An integral such as \(\int_{0}^{\infty} \tfrac{dx}{\sqrt{x}(1 + x)}\) is improper twice, at the endpoint \(0\) and along the infinite tail; split it at any interior point, say \(1\), and require each piece to converge by its own definition. In every case the procedure is the same: locate each point where the integral is improper, isolate it at one end of one integral, and apply the matching definition or the Comparison Theorem.

This section widened the integral without touching its foundations: an improper integral is a limit of proper ones, and \emph{convergent} is a verdict about that limit. Definitions \ref{def:8.1} and \ref{def:8.2} cover the infinite interval and the unbounded integrand, the \(p\)-test (Proposition \ref{prop:8.2}) marks the exact boundary between thick and thin tails, and the Comparison Theorem (Theorem \ref{thm:8.3}), proved with the monotone-sequence principle of Chapter 4, delivers verdicts when no antiderivative is available. These three results will be used again: Chapter 11’s Integral Test for infinite series is built directly on them. What this section does not do is produce numbers for integrals like \(\int_{1}^{\infty} e^{-x^{2}}\,dx\); it only certifies that the numbers exist. Producing them, to any stated accuracy, is the business of the chapter’s final section.


\sechead{8.7}{Approximate Integration}
The chapter has accumulated integrals that no formula reaches: \(\int_{1}^{\infty} e^{-x^{2}}\,dx\), certified finite in §8.6 but never evaluated; the perimeter of an ellipse, which §8.3 flagged; the arc lengths that §7.6 left standing as definite integrals. And in laboratory work the integrand is often not a formula at all, only a column of measured values. For all of these, the practical question is the same: how do we get the \emph{number}? The answer begins with an observation from §6.1’s first area estimates. For a decreasing integrand, a left-endpoint sum overshoots and a right-endpoint sum undershoots, so the true value sits between two numbers we can compute. Once the target is bracketed, the task is to tighten the bracket, and this section’s three rules tighten it dramatically: each replaces the integrand on every slice by something integrable exactly, first a constant, then a chord, then a parabola.

Fix the notation once. Divide \([a, b]\) into \(n\) slices of equal width \(h = \tfrac{b - a}{n}\), with nodes \(x_{i} = a + ih\) for \(i = 0, 1, \dots, n\), and write \(y_{i} = f(x_{i})\). For each rule \(R\), the \emph{error} is \(E_{R} = \int_{a}^{b} f(x)\,dx - R\): exact value minus approximation, so a positive error means the rule fell short.

\subsechead{The Midpoint and Trapezoidal Rules}
The first rule needs no new idea. Choose the midpoint \(\bar{x}_{i} = \tfrac{1}{2}(x_{i-1} + x_{i})\) of each slice and replace \(f\) there by the constant \(f(\bar{x}_{i})\):

\[ \int_{a}^{b} f(x)\,dx \;\approx\; M_{n} = h\,\bigl[f(\bar{x}_{1}) + f(\bar{x}_{2}) + \cdots + f(\bar{x}_{n})\bigr]. \tag{M} \]

This is a Riemann sum in the sense of Definition 6.2, with the sample points taken at midpoints, a choice that §6.2 reserved for exactly this use. For continuous \(f\), Theorem 6.1 guarantees that \(M_{n}\) converges to the integral as \(n\) grows; the question this section adds is how fast.

Unlike the two rules that follow, this rule never evaluates \(f\) at an endpoint of \([a, b]\), which is a practical advantage when the integrand is undefined or hard to evaluate there. The error bound below does not inherit that advantage. Theorem \ref{thm:8.4} requires \(f''\) to be continuous on all of \([a, b]\), endpoints included, so an integrand that is undefined or unbounded at an endpoint loses its guarantee even though the rule still returns a number.

The second rule replaces \(f\) on each slice by the chord joining \(\bigl(x_{i-1}, y_{i-1}\bigr)\) to \(\bigl(x_{i}, y_{i}\bigr)\). The region under a chord is a trapezoid with parallel sides \(y_{i-1}\) and \(y_{i}\) and width \(h\), so its area is \(h\cdot\tfrac{y_{i-1} + y_{i}}{2}\), and summing over the slices gives

\[ \int_{a}^{b} f(x)\,dx \;\approx\; T_{n} = \frac{h}{2}\,\bigl[y_{0} + 2y_{1} + 2y_{2} + \cdots + 2y_{n-1} + y_{n}\bigr]. \tag{T} \]

Every interior node is shared by two trapezoids, which is where the repeated factor \(2\) comes from. The same formula arrives by a second route. The left-endpoint and right-endpoint Riemann sums are \(L_{n} = h\sum_{i=1}^{n} y_{i-1}\) and \(R_{n} = h\sum_{i=1}^{n} y_{i}\), and averaging them termwise pairs each interior value with itself:

\[ \frac{L_{n} + R_{n}}{2} = \frac{h}{2}\sum_{i=1}^{n}\bigl(y_{i-1} + y_{i}\bigr) = \frac{h}{2}\bigl[y_{0} + 2y_{1} + \cdots + 2y_{n-1} + y_{n}\bigr] = T_{n}. \]

The averaging reading connects the rule to the bracketing picture above: for a monotone integrand the integral lies between \(L_{n}\) and \(R_{n}\), and \(T_{n}\) is the bracket’s midpoint, so the trapezoidal error is at most half the bracket’s width.

\begin{figureblock}
  \includegraphics[width=47.09mm]{ch08/figs/three-rules-panel-1}\hspace{3.86mm}\includegraphics[width=47.09mm]{ch08/figs/three-rules-panel-2}\hspace{3.86mm}\includegraphics[width=47.09mm]{ch08/figs/three-rules-panel-3}
\figcaption{fig:8.11}{The three rules on one integrand, \(y = 1/x\) on \([1, 3]\) with \(n = 4\): constants at the midpoints, chords at the nodes, and a parabola through each triple of nodes. The first triple’s arc is shaded; the open dots mark the second triple, whose arc lies too close to the curve to separate by eye, which is the point: each replacement follows the bending of the curve more closely than the last.}
\end{figureblock}
\begin{workedexample}
\begin{envexample}{Example}{ex:8.34}{}
Approximate \(\displaystyle\int_{1}^{2} \frac{dx}{x}\) with \(n = 5\) by the Trapezoidal Rule and by the Midpoint Rule, and compare each result with the true value.
\end{envexample}
\begin{envsolution}{Solution}{}{}
Here \(h = 0.2\), the nodes are \(1, 1.2, 1.4, 1.6, 1.8, 2\), and the midpoints are \(1.1, 1.3, 1.5, 1.7, 1.9\). The Trapezoidal Rule (T) gives

\[ T_{5} = \frac{0.2}{2}\Bigl[\frac{1}{1} + \frac{2}{1.2} + \frac{2}{1.4} + \frac{2}{1.6} + \frac{2}{1.8} + \frac{1}{2}\Bigr] \approx 0.695635, \]

and the Midpoint Rule (M) gives

\[ M_{5} = 0.2\,\Bigl[\frac{1}{1.1} + \frac{1}{1.3} + \frac{1}{1.5} + \frac{1}{1.7} + \frac{1}{1.9}\Bigr] \approx 0.691908. \]

The true value is \(\int_{1}^{2}\tfrac{dx}{x} = \ln 2 \approx 0.693147\), so the errors are

\[ E_{T} \approx -0.002488, \qquad E_{M} \approx +0.001239. \]

Two patterns are worth noticing, and both persist beyond this example. The midpoint error is about half the trapezoidal error in size, and the two have opposite signs: the chords of a curve that bends upward lie above it, so \(T_{5}\) overshoots, while the midpoint rectangles compensate through the tilt of the tangent, as Figure \ref{fig:8.12} shows. The error bounds of Theorem \ref{thm:8.4} below make the sizes unsurprising, since the midpoint bound has \(24\) in its denominator where the trapezoidal bound has \(12\), which makes the midpoint bound half the size. The bounds control only the errors’ magnitudes, not their signs or their exact ratio. \qedmark
\end{envsolution}
\end{workedexample}
\begin{figureblock}
  \includegraphics[width=90.48mm]{ch08/figs/midpoint-tangent-trap}
\figcaption{fig:8.12}{One slice of Example \ref{ex:8.34}’s integrand. Tilting the midpoint rectangle’s lid about the point of tangency onto the tangent line exchanges two equal triangles, so the rectangle and the tangent trapezoid hold the same area. On this upward-bending curve the tangent lies below the curve away from the midpoint while the chord lies above it, which is why the midpoint and trapezoidal errors pull in opposite directions.}
\end{figureblock}
\subsechead{Simpson’s Rule}
Chords match the curve’s value at two points but ignore its bending. The next upgrade replaces \(f\), on each \emph{pair} of adjacent slices, by the parabola through the three nodes the pair provides. One small lemma does all the work: the integral of a quadratic over a symmetric interval is determined by its three equally spaced values. For \(g(x) = Ax^{2} + Bx + C\) on \([-h, h]\),

\[ \int_{-h}^{h} g(x)\,dx = \frac{2Ah^{3}}{3} + 2Ch = \frac{h}{3}\bigl(2Ah^{2} + 6C\bigr), \]

while the values \(y_{0} = g(-h) = Ah^{2} - Bh + C\), \(y_{1} = g(0) = C\), \(y_{2} = g(h) = Ah^{2} + Bh + C\) combine to \(y_{0} + 4y_{1} + y_{2} = 2Ah^{2} + 6C\). Therefore

\[ \int_{-h}^{h} g(x)\,dx = \frac{h}{3}\bigl(y_{0} + 4y_{1} + y_{2}\bigr), \]

and \(B\) has dropped out entirely. Shifting the variable does not change an integral (Theorem 6.7 with a linear substitution), so the same weights \(1, 4, 1\) serve on any triple of equally spaced nodes. Now cover \([a, b]\) by \(\tfrac{n}{2}\) such triples, which requires \(n\) to be \emph{even}, and add: each pair of slices contributes \(\tfrac{h}{3}(y_{i-1} + 4y_{i} + y_{i+1})\) for \(i = 1, 3, \dots, n - 1\), an interior even-indexed node is shared by two pairs, and the sum of all the contributions is \emph{Simpson’s Rule}:

\[ \begin{aligned}
    \int_{a}^{b} f(x)\,dx \;\approx\; S_{n} = \frac{h}{3}\,\bigl[y_{0} &+ 4y_{1} + 2y_{2} + 4y_{3} + \cdots \\
      &+ 2y_{n-2} + 4y_{n-1} + y_{n}\bigr].
  \end{aligned} \tag{S} \]

\begin{envcaution}{Caution}{}{}
Simpson’s Rule pairs the slices, so it needs an even number of them. The weight pattern runs \(1, 4, 2, 4, 2, \dots, 2, 4, 1\): the \(4\)s sit on the odd-indexed nodes, the \(2\)s on the interior even-indexed ones, and only the two endpoints carry weight \(1\). An odd \(n\), or a misaligned weight pattern, silently produces a wrong answer.
\end{envcaution}
The rule carries the name of Thomas Simpson (1710–1761), although the same three-point weights were in use about a century before him.

\begin{workedexample}
\begin{envexample}{Example}{ex:8.35}{}
Approximate \(\displaystyle\int_{1}^{2} \frac{dx}{x}\) by Simpson’s Rule with \(n = 10\).
\end{envexample}
\begin{envsolution}{Solution}{}{}
Here \(h = 0.1\) and the nodes run \(1, 1.1, 1.2, \dots, 2\). Applying (S) with \(f(x) = 1/x\):

\[ \begin{aligned}
          S_{10} = \frac{0.1}{3}\Bigl[\frac{1}{1} &+ \frac{4}{1.1} + \frac{2}{1.2} + \frac{4}{1.3} + \frac{2}{1.4} + \frac{4}{1.5} \\
            &+ \frac{2}{1.6} + \frac{4}{1.7} + \frac{2}{1.8} + \frac{4}{1.9} + \frac{1}{2}\Bigr] \approx 0.693150.
        \end{aligned} \]

Against \(\ln 2 \approx 0.693147\), the error is about \(-3 \times 10^{-6}\). Example \ref{ex:8.34}’s rules, using half as many nodes, erred in the third decimal place; Simpson’s Rule with these ten slices is already correct to five. The jump matches the scales recorded in the bounds of Theorem \ref{thm:8.4} below: the trapezoidal bound shrinks like \(1/n^{2}\), the Simpson bound like \(1/n^{4}\), and the errors observed here behave accordingly. \qedmark
\end{envsolution}
\end{workedexample}
The size of that jump has a reason the lemma above already contains. Add a cubic term \(Dx^{3}\) to the quadratic \(g\). Over \([-h, h]\) it integrates to \(0\), because \(x^{3}\) is odd (Theorem 6.8), and in the weighted combination it contributes \(D(-h)^{3} + 4 \cdot 0 + Dh^{3} = 0\) as well. Neither side of the lemma changes, so the identity holds for cubics as it stands. Simpson’s Rule is therefore exact on every cubic polynomial, one degree beyond the parabolas it was built from, and that extra degree is what the \emph{fourth} derivative in the next theorem records.

\subsechead{Error bounds}
An estimate is useful only when its error is controlled. The bending of \(f\) is what chords and rectangles cannot follow, so it is the size of the second derivative, and for Simpson’s Rule the fourth, that governs the error. Bounds for all three rules follow; each \(K\) exists whenever the named derivative is continuous, by the Extreme Value Theorem (Theorem 4.9(a)), and any valid bound serves in place of the exact maximum.

\begin{envtheorem}{Theorem}{thm:8.4}{Error Bounds for Numerical Integration}
Suppose \(f''\) is continuous on \([a, b]\) and \(\lvert f''(x)\rvert \le K\) for all \(x\) in \([a, b]\). Then

\[ \lvert E_{T}\rvert \le \frac{K\,(b - a)^{3}}{12\,n^{2}} \qquad\text{and}\qquad \lvert E_{M}\rvert \le \frac{K\,(b - a)^{3}}{24\,n^{2}}. \]

If moreover \(f^{(4)}\) is continuous on \([a, b]\) and \(\lvert f^{(4)}(x)\rvert \le K_{4}\) there, then for even \(n\)

\[ \lvert E_{S}\rvert \le \frac{K_{4}\,(b - a)^{5}}{180\,n^{4}}. \]
\end{envtheorem}
\begin{envcaution}{Caution}{}{}
Theorem \ref{thm:8.4} is stated without proof. The proofs measure how far a function strays from its chords and parabolas, and the right instrument for that is Taylor’s theorem with remainder, which Chapter 11 establishes; the bounds above are proved there. This is the second and last of the results this chapter takes on credit; the other is the factorization fact of §8.4.
\end{envcaution}
The theorem confirms the exactness noted above from the other direction: a cubic has \(f^{(4)} = 0\), so \(K_{4} = 0\) is a valid bound and the Simpson bound reads \(\lvert E_{S}\rvert \le 0\).

\begin{workedexample}
\begin{envexample}{Example}{ex:8.36}{}
How large must \(n\) be to guarantee that the Trapezoidal Rule approximates \(\displaystyle\int_{1}^{2}\frac{dx}{x}\) with error at most \(10^{-4}\)? How large for the Midpoint Rule?
\end{envexample}
\begin{envsolution}{Solution}{}{}
For \(f(x) = 1/x\) we have \(f''(x) = 2/x^{3}\), which is positive and decreasing on \([1, 2]\), so \(\lvert f''(x)\rvert \le f''(1) = 2\), and \(K = 2\) serves. With \(b - a = 1\), Theorem \ref{thm:8.4} gives

\[ \lvert E_{T}\rvert \le \frac{2}{12 n^{2}} = \frac{1}{6 n^{2}}, \]

and \(\tfrac{1}{6n^{2}} \le 10^{-4}\) holds exactly when \(n^{2} \ge \tfrac{10^{4}}{6}\), that is, \(n \ge \sqrt{10^{4}/6} \approx 40.82\). Since \(n\) is a whole number and the requirement is a guarantee, round \emph{up}: \(n = 41\) suffices. For the Midpoint Rule the same computation with the constant \(24\) gives \(n^{2} \ge \tfrac{10^{4}}{12}\), so \(n \ge \sqrt{10^{4}/12} \approx 28.87\), and \(n = 29\) suffices. Rounding down would break the guarantee: the bound at \(n = 40\) exceeds the tolerance. The bounds are worst-case: Example \ref{ex:8.34}’s actual errors were smaller than Theorem \ref{thm:8.4} predicts, and that is the usual experience. \qedmark
\end{envsolution}
\end{workedexample}
\begin{workedexample}
\begin{envexample}{Example}{ex:8.37}{}
Approximate \(\displaystyle\int_{0}^{1} e^{-x^{2}}\,dx\) by Simpson’s Rule with \(n = 10\), and state how many decimal places the result is guaranteed to give.
\end{envexample}
\begin{envsolution}{Solution}{}{}
The integrand is continuous on \([0, 1]\), so the integral exists (Theorem 6.1). What is missing is an antiderivative to hand to the Fundamental Theorem (Theorem 6.4): none of this book’s techniques produces one, which is why §8.5 set this integrand aside. The rules of this section do not need one.

With \(n = 10\) the width is \(h = 0.1\), the nodes are \(0, 0.1, 0.2, \dots, 1\), and \(10\) is even, so (S) applies:

\[ \begin{aligned}
          S_{10} = \frac{0.1}{3}\Bigl[e^{0} &+ 4e^{-0.01} + 2e^{-0.04} + \cdots \\
            &+ 2e^{-0.64} + 4e^{-0.81} + e^{-1}\Bigr] \approx 0.746825.
        \end{aligned} \]

For the bound, differentiate four times. With \(f(x) = e^{-x^{2}}\), the chain rule (Theorem 3.3), the derivative of the exponential (Theorem 4.8), and the product rule (Theorem 2.6) give

\[ f''(x) = \bigl(4x^{2} - 2\bigr)e^{-x^{2}}, \qquad f^{(4)}(x) = \bigl(16x^{4} - 48x^{2} + 12\bigr)e^{-x^{2}}. \]

On \([0, 1]\) we have \(0 \le x^{2} \le 1\) and \(0 < e^{-x^{2}} \le 1\). Bounding each term of the factor \(16x^{4} - 48x^{2} + 12\) separately, we find that it lies between \(0 - 48 + 12 = -36\) and \(16 - 0 + 12 = 28\). Since \(0 < e^{-x^{2}} \le 1\), multiplying by \(e^{-x^{2}}\) cannot increase its absolute value, so \(\lvert f^{(4)}(x)\rvert \le 36\) throughout. Take \(K_{4} = 36\). The theorem requires a valid bound, not the exact maximum; a generous bound weakens the guarantee without invalidating it, and here it saves the work of locating the maximum exactly.

Theorem \ref{thm:8.4} now gives, with \(b - a = 1\) and \(n = 10\),

\[ \lvert E_{S}\rvert \le \frac{36 \cdot 1^{5}}{180 \cdot 10^{4}} = 2 \times 10^{-5}. \]

The true value therefore lies between \(0.746805\) and \(0.746845\). Both ends round to \(0.7468\), so

\[ \int_{0}^{1} e^{-x^{2}}\,dx = 0.7468 \]

to four decimal places. Examples \ref{ex:8.34} and \ref{ex:8.35} measured their errors against a known value, \(\ln 2\). Here no such value is available, and the bound is the only statement of accuracy that can be made. It is still enough to name four decimal places, which is what the question asked for. \qedmark
\end{envsolution}
\end{workedexample}
\begin{workedexample}
\begin{envexample}{Example}{ex:8.38}{}
A car’s speedometer is read every \(30\) seconds during a three-minute drive, giving speeds, in meters per second, of

\[ 8, \quad 10, \quad 12, \quad 11, \quad 13, \quad 14, \quad 12 \]

at \(t = 0, 30, 60, 90, 120, 150, 180\). Estimate the distance traveled.
\end{envexample}
\begin{envsolution}{Solution}{}{}
Distance is the integral of speed (the Net Change Theorem, Theorem 6.5), so the target is \(\int_{0}^{180} v(t)\,dt\) with \(v\) known only at seven equally spaced times. Seven values mean \(n = 6\) slices of width \(h = 30\), and \(6\) is even, so Simpson’s Rule applies:

\[ \begin{aligned}
          S_{6} &= \frac{30}{3}\Bigl[8 + 4(10) + 2(12) + 4(11) + 2(13) + 4(14) + 12\Bigr] \\
            &= 10 \cdot 210 = 2100.
        \end{aligned} \]

The estimated distance is \(2100\) meters, an average speed of about \(11.7\) m/s, or \(42\) km/h, consistent with the readings. Two honest notes are in order. There is no error bound to quote: Theorem \ref{thm:8.4} needs derivative information, and a data column carries none, so the estimate’s quality rests on the smoothness of the actual speed between readings. And the rule required nothing but arithmetic on the data, which is why numerical integration, not antidifferentiation, is what measurement-driven work uses. \qedmark
\end{envsolution}
\end{workedexample}
\begin{envcaution}{Caution}{}{}
All three rules, and Theorem \ref{thm:8.4} with them, are built on equally spaced nodes, that is, on slices of equal width \(h = \tfrac{b - a}{n}\). Measured data does not always arrive that way. Readings at unequal spacings can still be handled by adding trapezoid areas one piece at a time, since each piece is a trapezoid whatever its width. The weight pattern of (S) and the bounds of Theorem \ref{thm:8.4}, however, are not available in that setting.
\end{envcaution}
\begin{figureblock}
  \includegraphics[width=126.5mm]{ch08/figs/simpson-data-speed}
\figcaption{fig:8.13}{Example \ref{ex:8.38}’s data: seven speedometer readings and the three parabolic arcs that Simpson’s Rule fits through consecutive triples, meeting at the shared nodes \(t = 60\) and \(t = 120\). No curve is drawn between the readings, because none is known; the first triple happens to be collinear, so its arc is a straight segment.}
\end{figureblock}
\begin{envcaution}{Caution}{}{}
One of the integrals named at the start of this section still needs a comment. All three rules act on a finite interval, so \(\int_{1}^{\infty} e^{-x^{2}}\,dx\) cannot be handed to them as it stands. Split it instead: \(\int_{1}^{\infty} = \int_{1}^{t} + \int_{t}^{\infty}\), approximate the first piece by any of the rules, and bound the discarded tail. Example \ref{ex:8.32} supplies that bound already: its comparison \(e^{-x^{2}} \le e^{-x}\) holds for every \(x \ge 1\). Comparing on \([t, \infty)\), where \(\int_{t}^{\infty} e^{-x}\,dx = e^{-t}\), gives \(\int_{t}^{\infty} e^{-x^{2}}\,dx \le e^{-t}\) (Theorem \ref{thm:8.3}(a)). The total error is then at most the rule’s error on \([1, t]\) plus that tail bound, and \(t\) is chosen large enough to make the tail as small as the target accuracy requires.
\end{envcaution}
\subsechead{Chapter summary}
This chapter set out to finish what the Fundamental Theorem started: if evaluation reduces to antidifferentiation, then antidifferentiation needs a toolkit. Section 8.1 built the first tool from the product rule, integration by parts, \(\int u\,dv = uv - \int v\,du\) (Theorems \ref{thm:8.1} and \ref{thm:8.2}), with a strategy for choosing the parts (Strategy \ref{strat:8.1}) and, as a by-product, the settlement of Chapter 7’s shell-versus-washer question in the monotone case. Section 8.2 turned trigonometric identities into substitutions, filled the antiderivative table of §6.4 with \(\int\tan x\,dx = \ln\lvert\sec x\rvert + C\) and its companions (Proposition \ref{prop:8.1}), and recorded \(\int\sec^{3}x\,dx\) for later use. Section 8.3 reparametrized quadratic radicals away with the three restricted substitutions of Strategy \ref{strat:8.4}, and §8.4 sent every rational function through the pipeline of Strategy \ref{strat:8.5}, dividing, factoring on the credit of the Fundamental Theorem of Algebra, decomposing by Proposition A.7, and integrating fragment by fragment, with \(\int\frac{dx}{x^{2}+a^{2}} = \frac{1}{a}\arctan\frac{x}{a} + C\) as the recurring final step. Section 8.5 compressed the toolkit into a four-step decision procedure (Strategy \ref{strat:8.6}) and marked its boundary: some integrands, such as \(e^{-x^{2}}\), cannot be handled by any of it.

The last two sections changed the question. Section 8.6 extended the integral itself to infinite intervals and unbounded integrands (Definitions \ref{def:8.1} and \ref{def:8.2}), proved the \(p\)-test boundary \(\int_{1}^{\infty} x^{-p}\,dx\) convergent exactly for \(p > 1\) (Proposition \ref{prop:8.2}), and proved the Comparison Theorem (Theorem \ref{thm:8.3}) with Chapter 4’s monotone-sequence principle, so that convergence can be certified without evaluation. Section 8.7 then produced the numbers: the Midpoint, Trapezoidal, and Simpson’s Rules (M), (T), (S), the last derived from the three-point parabola lemma, with error bounds (Theorem \ref{thm:8.4}) that let us choose \(n\) in advance. The chapter closes with exactly two results taken on credit, both named where they were incurred: the real-factorization fact behind §8.4, borrowed permanently from algebra, and the proofs of Theorem \ref{thm:8.4}, which Chapter 11’s Taylor theorem delivers. Chapter 11 also draws on this chapter directly: its Integral Test is built on Definition \ref{def:8.1}, Proposition \ref{prop:8.2}, and Theorem \ref{thm:8.3}. Before Chapter 11’s infinite series arrive, Chapter 9 puts the completed toolkit to work on equations whose unknowns are functions, where every solution method begins by asking the questions this chapter taught: what form is this, and which technique fits?
\end{document}
